\chapter{Convergence}
\label{chap:pet-convergence}

\newcommand{\Ric}{\operatorname{Ric}}
\newcommand{\Hess}{\operatorname{Hess}}
\newcommand{\vol}{\operatorname{vol}}
\newcommand{\diam}{\operatorname{diam}}
\newcommand{\inj}{\operatorname{inj}}
\newcommand{\sec}{\operatorname{sec}}
\newcommand{\Cap}{\operatorname{Cap}}
\newcommand{\Cov}{\operatorname{Cov}}
\newcommand{\dGH}{d_{G\text{-}H}}

Faithful transcription of Petersen, \emph{Riemannian Geometry} (GTM 171, 3rd
ed.), Chapter 11. The chapter opens with a motivational overview (no numbered
statements) of convergence notions for Riemannian manifolds, then develops:
the Hausdorff and Gromov--Hausdorff distances on compact or proper metric
spaces and the resulting metric space $(\mathcal{M},\dGH)$ of compact metric
spaces, together with pointed convergence, convergence of maps, and Gromov's
precompactness criteria (\S11.1); the Hölder spaces $C^{m,\alpha}$ and the
Schauder elliptic estimates used to construct harmonic coordinates (\S11.2);
the $C^{m,\alpha}$ and harmonic norms of a pointed Riemannian manifold, the
induced notion of pointed $C^{m,\alpha}$ convergence, and Cheeger's
Fundamental Theorem of Convergence Theory together with its finiteness
corollaries (\S11.3); and the geometric applications of the theory —
Anderson's harmonic-coordinate control of the metric by Ricci curvature, the
Convergence Theorem of Riemannian Geometry, Cheeger's injectivity-radius
lemma, and the Ricci/sectional curvature pinching theorems (\S11.4) — closing
with Petersen's bibliographic remarks on further study of convergence and
collapsing (\S11.5) and his exercise set (\S11.6).

\section{Gromov-Hausdorff Convergence}

\subsection{Hausdorff Versus Gromov Convergence}

\begin{definition}[Hausdorff distance]
  \label{def:pet-ch11-hausdorff-distance}
  \lean{PetersenLib.hausdorffDistance}
  For a metric space $(X,|\cdot|)$ and $A,B\subset X$, set
  $d(A,B)=\inf\{|ab|\mid a\in A,\,b\in B\}$,
  $B(A,\varepsilon)=\{x\in X\mid|xA|<\varepsilon\}$ (writing $|xA|=d(\{x\},A)$),
  and define the \emph{Hausdorff distance}
  \[
    d_H(A,B)=\inf\{\varepsilon\mid A\subset B(B,\varepsilon),\,B\subset B(A,\varepsilon)\}.
  \]
  Thus $d(A,B)$ is small when some points of $A,B$ are close, while
  $d_H(A,B)$ is small exactly when every point of $A$ is close to a point of
  $B$ and vice versa. The Hausdorff distance defines a metric on the compact
  subsets of $X$, and this collection is compact when $X$ is compact.
  \source{petersen:page-0409}
\end{definition}

\begin{convention}[compact and proper metric spaces]
  \label{conv:pet-ch11-compact-proper-spaces}
  \lean{PetersenLib.ProperMetricSpace}
  Throughout, only compact or \emph{proper} metric spaces are considered:
  spaces whose distance function is proper, i.e.\ all closed balls are
  compact. Proper metric spaces are automatically separable, complete, and
  locally compact.
  \source{petersen:page-0409}
\end{convention}

\begin{definition}[admissible metric, Gromov--Hausdorff distance]
  \label{def:pet-ch11-gromov-hausdorff-distance}
  \lean{PetersenLib.admissibleMetric,PetersenLib.gromovHausdorffDistance}
  \uses{def:pet-ch11-hausdorff-distance}
  For metric spaces $X,Y$, an \emph{admissible metric} on the disjoint union
  $X\cup Y$ is a metric that extends the given metrics on $X$ and on $Y$. The
  \emph{Gromov--Hausdorff distance} is
  \[
    \dGH(X,Y)=\inf\{d_H(X,Y)\mid\text{admissible metrics on }X\cup Y\}:
  \]
  one places a metric on $X\cup Y$ extending those on $X,Y$ so as to make $X$
  and $Y$ as Hausdorff-close as possible — equivalently, one tries to define
  distances between points of $X$ and $Y$ without violating the triangle
  inequality.
  \source{petersen:page-0409}
\end{definition}

\begin{example}[Example 11.1.1 — distance to a point]
  \label{ex:pet-ch11-gh-distance-radius}
  \lean{PetersenLib.ghDistance_le_radius}
  \uses{def:pet-ch11-gromov-hausdorff-distance}
  If $Y$ is a one-point space, then
  $\dGH(X,Y)\le\operatorname{rad}X=\inf_{y}\sup_{x\in X}|xy|$, the radius of
  the smallest ball covering $X$.
  \source{petersen:page-0409}
\end{example}

\begin{example}[Example 11.1.2 — diameter bound]
  \label{ex:pet-ch11-gh-distance-diameter-bound}
  \lean{PetersenLib.ghDistance_le_halfDiam}
  \uses{def:pet-ch11-gromov-hausdorff-distance}
  If $\diam X,\diam Y\le D$, using $|xy|=D/2$ for all $x\in X$, $y\in Y$ as an
  admissible metric on $X\cup Y$ shows $\dGH(X,Y)\le D/2$.
  \source{petersen:page-0409}
\end{example}

\begin{proposition}[Prop.\ 11.1.3 — distance zero implies isometric]
  \label{prop:pet-ch11-gh-distance-zero-isometric}
  \lean{PetersenLib.ghDistance_eq_zero_iff_isometric}
  \uses{def:pet-ch11-gromov-hausdorff-distance}
  If $X$ and $Y$ are compact metric spaces with $\dGH(X,Y)=0$, then $X$ and
  $Y$ are isometric.
  \source{petersen:page-0410}
\end{proposition}
\begin{proof}
  \uses{prop:pet-ch11-gh-distance-zero-isometric}
  Choose metrics $|\cdot|_i$ on $X\cup Y$ with Hausdorff distance between $X$
  and $Y$ $<i^{-1}$. Then find (possibly discontinuous) maps
  $I_i:X\to Y$ with $|xI_i(x)|_i\le i^{-1}$ and $J_i:Y\to X$ with
  $|yJ_i(y)|_i\le i^{-1}$. The triangle inequality and the fact that
  $|\cdot|_i$ restricted to $X$ (resp.\ $Y$) is the given metric yield
  \[
    |I_i(x_1)I_i(x_2)|\le2i^{-1}+|x_1x_2|,\qquad
    |J_i(y_1)J_i(y_2)|\le2i^{-1}+|y_1y_2|,
  \]
  \[
    |xJ_i\circ I_i(x)|\le2i^{-1},\qquad |yI_i\circ J_i(y)|\le2i^{-1}.
  \]
  Construct $I:X\to Y$, $J:Y\to X$ as limits of the $I_i,J_i$ as in the proof
  of the Arzelà--Ascoli lemma: for each $x$ the sequence $(I_i(x))$ has an
  accumulation point since $Y$ is compact; fixing a countable dense
  $A\subset X$, a diagonal argument gives a subsequence $I_{i_j}$ with
  $I_{i_j}(a)\to I(a)$ for all $a\in A$. The first inequality shows $I$ is
  distance decreasing on $A$, hence uniformly continuous, and extends
  uniquely to a distance-decreasing map $I:X\to Y$; likewise for $J:Y\to X$.
  The last two inequalities pass to the limit to show $I,J$ are mutual
  inverses, so both are isometries.
\end{proof}

\begin{remark}[$(\mathcal{M},\dGH)$ as a metric space]
  \label{rem:pet-ch11-gh-pseudo-metric-quotient}
  \lean{PetersenLib.compactMetricSpaceQuotient}
  \uses{prop:pet-ch11-gh-distance-zero-isometric}
  Symmetry and the triangle inequality for $\dGH$ are straightforward, so
  $\dGH$ is a pseudo-metric on the collection $\mathcal{M}$ of compact metric
  spaces; by \cref{prop:pet-ch11-gh-distance-zero-isometric} the equivalence
  classes (up to isometry) form a genuine metric space $(\mathcal{M},\dGH)$.
  \source{petersen:page-0410}
\end{remark}

\begin{example}[Example 11.1.4 — $\varepsilon$-dense finite approximation]
  \label{ex:pet-ch11-epsilon-dense-approximation}
  \lean{PetersenLib.epsilonDense,PetersenLib.ghDistance_le_epsilon_of_epsilonDense}
  \uses{def:pet-ch11-gromov-hausdorff-distance,def:pet-ch11-hausdorff-distance}
  Let $X$ be compact and $A\subset X$ finite with $d_H(A,X)\le\varepsilon$
  (metric on $A$ induced from $X$); such $A$ is \emph{$\varepsilon$-dense} in
  $X$. Then $\dGH(X,A)\le\varepsilon$; since $X$ is compact, finite
  $\varepsilon$-dense subsets exist for every $\varepsilon>0$.
  \source{petersen:page-0410,petersen:page-0411}
\end{example}
\begin{proof}
  \uses{ex:pet-ch11-epsilon-dense-approximation}
  For small $\delta>0$ put $|xa|_{X\cup A}=\delta+|xa|_X$ for $x\in X$,
  $a\in A$; this is an admissible metric on $X\cup A$ for which
  $d_H(X,A)<\varepsilon+\delta$ (using $d_H(A,X)\le\varepsilon$), so
  $\dGH(X,A)<\varepsilon+\delta$. Letting $\delta\to0$ gives the estimate.
\end{proof}

\begin{example}[Example 11.1.5 — correspondences bound $\dGH$ by $3\varepsilon$]
  \label{ex:pet-ch11-finite-approximation-3eps}
  \lean{PetersenLib.ghDistance_le_threeEpsilon_of_correspondence}
  \uses{ex:pet-ch11-epsilon-dense-approximation}
  Suppose $A=\{x_1,\dots,x_k\}\subset X$, $B=\{y_1,\dots,y_k\}\subset Y$ are
  $\varepsilon$-dense subsets with
  $\big||x_ix_j|-|y_iy_j|\big|\le\varepsilon$ for $1\le i,j\le k$. Then
  $\dGH(X,Y)\le3\varepsilon$.
  \source{petersen:page-0411,petersen:page-0412}
\end{example}
\begin{proof}
  \uses{ex:pet-ch11-finite-approximation-3eps}
  By \cref{ex:pet-ch11-epsilon-dense-approximation} and the triangle
  inequality it suffices to show $\dGH(A,B)\le\varepsilon$, i.e.\ to exhibit a
  metric on $A\cup B$ making $A,B$ $\varepsilon$-Hausdorff close. Define
  $|x_iy_i|=\varepsilon$ and
  $|x_iy_j|=\min_k\{|x_ix_k|+\varepsilon+|y_jy_k|\}$; this extends the given
  metrics on $A,B$, identifies no points across $A,B$, and is symmetric. For
  the triangle inequality it suffices to check
  $|x_iy_j|\le|x_iz|+|y_jz|$ and $|x_ix_j|\le|y_kx_i|+|y_kx_j|$ (the case
  $|y_iy_j|\le|x_ky_i|+|x_ky_j|$ is symmetric). For the first, take
  $z=x_k$ and choose $l$ with
  $|y_jx_k|=\varepsilon+|y_jy_l|+|x_lx_k|$; then
  \[
    |x_ix_k|+|y_jx_k|=|x_ix_k|+\varepsilon+|y_jy_l|+|x_lx_k|
    \ge|x_iy_l|+\varepsilon+|y_jy_l|\ge|x_iy_j|.
  \]
  For the second, choose $l,m$ with
  $|y_kx_i|=|y_ky_l|+\varepsilon+|x_lx_i|$ and
  $|y_kx_j|=|y_ky_m|+\varepsilon+|x_mx_j|$; the correspondence hypothesis
  $\big||x_ix_j|-|y_iy_j|\big|\le\varepsilon$ then gives
  \[
    |y_kx_i|+|y_kx_j|
    =|y_ky_l|+\varepsilon+|x_lx_i|+|y_ky_m|+\varepsilon+|x_mx_j|
    \ge|x_kx_l|+|x_lx_i|+|x_kx_m|+|x_mx_j|\ge|x_ix_j|,
  \]
  using the triangle inequality in $X$ twice on the last step.
\end{proof}

\begin{example}[Example 11.1.6 — lens spaces collapse to $S^2(1/2)$]
  \label{ex:pet-ch11-lens-spaces-limit}
  \lean{PetersenLib.lensSpace_ghConverges_toSphere}
  \uses{ex:pet-ch11-finite-approximation-3eps}
  Let $M_k=S^3/\mathbb{Z}_k$ with the metric induced from $S^3(1)$. The
  Riemannian submersion $M_k\to S^2(1/2)$ has fibres of diameter
  $2\pi/k\to0$; using \cref{ex:pet-ch11-finite-approximation-3eps} it follows
  that $M_k\to S^2(1/2)$ in the Gromov--Hausdorff topology.
  \source{petersen:page-0412}
\end{example}

\begin{example}[Example 11.1.7 — Berger spheres collapse to $S^2(1/2)$]
  \label{ex:pet-ch11-berger-spheres-limit}
  \lean{PetersenLib.bergerSphere_ghConverges_toSphere}
  \uses{ex:pet-ch11-lens-spaces-limit}
  Likewise the Berger metrics $(S^3,g_\varepsilon)\to S^2(1/2)$ as
  $\varepsilon\to0$. In both examples the volume goes to zero while the
  curvatures and diameters stay uniformly bounded; in the Berger case the
  manifolds are even simply connected, yet the topology changes drastically —
  in the lens-space case the fundamental groups of the sequence elements are
  even mutually different.
  \source{petersen:page-0412}
\end{example}

\begin{proposition}[Prop.\ 11.1.8 — $(\mathcal{M},\dGH)$ separable and complete]
  \label{prop:pet-ch11-gh-space-separable-complete}
  \lean{PetersenLib.gromovHausdorffSpace_separable_complete}
  \uses{rem:pet-ch11-gh-pseudo-metric-quotient}
  The metric space $(\mathcal{M},\dGH)$ is separable and complete.
  \source{petersen:page-0412,petersen:page-0413}
\end{proposition}
\begin{proof}
  \uses{prop:pet-ch11-gh-space-separable-complete}
  \emph{Separability.} Finite metric spaces are dense in $\mathcal{M}$ by
  \cref{ex:pet-ch11-epsilon-dense-approximation}, so it suffices to note that
  the countable collection of finite metric spaces all of whose pairwise
  distances are rational is dense as well.

  \emph{Completeness.} Let $\{X_n\}$ be Cauchy; it suffices to produce a
  convergent subsequence. Choose $\{X_i\}$ with
  $\dGH(X_i,X_{i+1})<2^{-i}$ and metrics $|\cdot|_{i,i+1}$ on $X_i\cup X_{i+1}$
  making them $2^{-i}$-Hausdorff close. Define $|\cdot|_{i,i+j}$ on
  $X_i\cup X_{i+j}$ by
  \[
    |x_i\,x_{i+j}|_{i,i+j}
    =\min_{\{x_{i+k}\in X_{i+k}\}}\Big\{\sum_{k=0}^{j-1}|x_{i+k}x_{i+k+1}|\Big\},
  \]
  giving a metric on $Y=\bigcup_iX_i$ with $d_H(X_i,X_{i+j})\le2^{-i+1}$. Let
  \[
    \hat X=\big\{\{x_i\}\mid x_i\in X_i,\ |x_iy_j|\to0\text{ as }i,j\to\infty\big\},
  \]
  with pseudo-metric $|\{x_i\}\{y_i\}|=\lim_i|x_iy_i|$; the quotient by
  $\{x_i\}\sim\{y_i\}\iff|\{x_i\}\{y_i\}|=0$ is a metric space $X$. Extend the
  metric on $Y$ to $X\cup Y$ by
  $|x\,\{x_i\}|=\lim_i|x_ix_i|$. Since $d_H(X_i,X_{i+1})\le2^{-i}$, for any
  $x_i\in X_i$ one can build $\{x_{i+j}\}\in\hat X$ with $x_{i+0}=x_i$ and
  $|x_{i+j}x_{i+j+1}|\le2^{-j}$, so $|x_i\{x_{i+j}\}|\le2^{-i+1}$: every
  $X_i$ is $2^{-i+1}$-close to $X$. Conversely any sequence $\{x_i\}$ is
  equivalent to one $\{y_i\}$ with $|y_k\{y_i\}|\le2^{-k+1}$, so $X$ is
  $2^{-i+1}$-close to $X_i$. Hence $X_i\to X$ in $\dGH$.
\end{proof}

\subsection{Pointed Convergence}

\begin{definition}[pointed Gromov--Hausdorff distance]
  \label{def:pet-ch11-pointed-gh-distance}
  \lean{PetersenLib.pointedGromovHausdorffDistance}
  \uses{def:pet-ch11-gromov-hausdorff-distance,def:pet-ch11-hausdorff-distance}
  For pointed metric spaces $(X,x)$, $(Y,y)$, the \emph{pointed
  Gromov--Hausdorff distance} is
  \[
    \dGH\big((X,x),(Y,y)\big)=\inf\{d_H(X,Y)+|xy|\},
  \]
  the infimum over admissible metrics on $X\cup Y$ of the Hausdorff distance
  plus the distance between the marked points.
  \source{petersen:page-0414}
\end{definition}

\begin{definition}[pointed Gromov--Hausdorff topology]
  \label{def:pet-ch11-pointed-gh-topology}
  \lean{PetersenLib.PointedGromovHausdorffConvergence}
  \uses{def:pet-ch11-pointed-gh-distance}
  On the collection $\mathcal{M}_*=\{(X,x,|\cdot|)\}$ of proper pointed
  metric spaces, $(X_i,x_i,|\cdot|_i)\to(X,x,|\cdot|)$ in the \emph{pointed
  Gromov--Hausdorff topology} if for every $R$ there is a sequence $R_i\to R$
  such that the closed balls
  $\big(\overline{B}(x_i,R_i),x_i,|\cdot|_i\big)\to\big(\overline{B}(x,R),x,|\cdot|\big)$
  converge in the pointed Gromov--Hausdorff metric of
  \cref{def:pet-ch11-pointed-gh-distance}.
  \source{petersen:page-0414}
\end{definition}

\subsection{Convergence of Maps}

\begin{definition}[convergence of maps]
  \label{def:pet-ch11-convergence-of-maps}
  \lean{PetersenLib.MapsConvergeTo}
  \uses{def:pet-ch11-pointed-gh-topology}
  Given $f_k:X_k\to Y_k$ with $X_k\to X$ and $Y_k\to Y$ (Gromov--Hausdorff),
  $f_k$ \emph{converges} to $f:X\to Y$ if for every sequence $x_k\in X_k$
  with $x_k\to x\in X$ one has $f_k(x_k)\to f(x)$. This depends on the choice
  of Hausdorff-convergent ambient metrics. It is a strong notion: if
  $X_k=X$, $Y_k=Y$, $f_k=f$ for all $k$, then $f$ converges to itself only if
  $f$ is continuous.
  \source{petersen:page-0414}
\end{definition}

\begin{remark}[convergence of maps as an extension problem]
  \label{rem:pet-ch11-map-convergence-as-extension}
  \lean{PetersenLib.mapConvergence_iff_continuousExtension}
  \uses{def:pet-ch11-convergence-of-maps}
  Convergence of maps preserves being distance preserving or a submetry. A
  sequence $f_k:X_k\to Y_k$ can be regarded as one continuous map
  $F:\bigcup_iX_i\to Y\cup\bigcup_iY_i$; the sequence converges iff $F$ has a
  continuous extension $X\cup\bigcup_iX_i\to Y\cup\bigcup_iY_i$ restricting to
  the limit map on $X$. When the $X_i$ are compact this is equivalent to $F$
  being uniformly continuous.
  \source{petersen:page-0415}
\end{remark}

\begin{definition}[equicontinuous family]
  \label{def:pet-ch11-equicontinuous-family}
  \lean{PetersenLib.EquicontinuousFamily}
  \uses{def:pet-ch11-convergence-of-maps}
  A family $f_k:X_k\to Y_k$ is \emph{equicontinuous} if for every
  $\varepsilon>0$ and $x_k\in X_k$ there is $\delta>0$ with
  $f_k(B(x_k,\delta))\subset B(f_k(x_k),\varepsilon)$ for all $k$. Uniformly
  Lipschitz families are equicontinuous.
  \source{petersen:page-0415}
\end{definition}

\begin{lemma}[Lemma 11.1.9 — Gromov--Hausdorff Arzelà--Ascoli]
  \label{lem:pet-ch11-arzela-ascoli-gh}
  \lean{PetersenLib.arzelaAscoli_gromovHausdorff}
  \uses{def:pet-ch11-equicontinuous-family}
  An equicontinuous family $f_k:X_k\to Y_k$, with $X_k\to X$ and $Y_k\to Y$
  in the (pointed) Gromov--Hausdorff topology, has a convergent subsequence.
  In the pointed case, one also assumes each $f_k$ preserves the base point.
  \source{petersen:page-0415}
\end{lemma}
\begin{proof}
  \uses{lem:pet-ch11-arzela-ascoli-gh}
  Choose dense subsets $A_i=\{a_1^i,a_2^i,\dots\}\subset X_i$ with
  $a_j^i\to a_j\in X$ as $i\to\infty$, so $A=\{a_j\}\subset X$ is dense. A
  diagonal argument, using equicontinuity, produces a subsequence of the
  $f_k$ converging on the sequences $a_j^i$; equicontinuity then upgrades
  this to convergence of the maps in the sense of
  \cref{def:pet-ch11-convergence-of-maps}.
\end{proof}

\subsection{Compactness of Classes of Metric Spaces}

\begin{definition}[capacity and covering functions]
  \label{def:pet-ch11-capacity-covering-functions}
  \lean{PetersenLib.capacityFunction,PetersenLib.coveringFunction}
  For a compact metric space $X$, let
  $\Cap(\varepsilon)=\Cap_X(\varepsilon)$ be the maximum number of disjoint
  $\varepsilon/2$-balls in $X$, and $\Cov(\varepsilon)=\Cov_X(\varepsilon)$
  the minimum number of $\varepsilon$-balls needed to cover $X$.
  \source{petersen:page-0416}
\end{definition}

\begin{lemma}[covering is bounded by capacity]
  \label{lem:pet-ch11-cov-le-cap}
  \lean{PetersenLib.coveringFunction_le_capacityFunction}
  \uses{def:pet-ch11-capacity-covering-functions}
  $\Cov(\varepsilon)\le\Cap(\varepsilon)$.
  \source{petersen:page-0416}
\end{lemma}
\begin{proof}
  \uses{lem:pet-ch11-cov-le-cap}
  Select a maximal disjoint family $B(x_i,\varepsilon/2)$ and consider the
  balls $B(x_i,\varepsilon)$. If these do not cover $X$, some
  $x\in X-\bigcup B(x_i,\varepsilon)$ gives $B(x,\varepsilon/2)$ disjoint from
  every $B(x_i,\varepsilon/2)$, contradicting maximality of the disjoint
  family; hence the $B(x_i,\varepsilon)$ cover $X$, and there are
  $\Cap(\varepsilon)$ of them.
\end{proof}

\begin{lemma}[covering/capacity under Gromov--Hausdorff proximity]
  \label{lem:pet-ch11-gh-close-cov-cap-comparison}
  \lean{PetersenLib.coveringFunction_capacityFunction_ghClose}
  \uses{def:pet-ch11-capacity-covering-functions,def:pet-ch11-gromov-hausdorff-distance}
  If $\dGH(X,Y)<\delta$, then
  $\Cov_X(\varepsilon+2\delta)\le\Cov_Y(\varepsilon)$ and
  $\Cap_X(\varepsilon)\ge\Cap_Y(\varepsilon+2\delta)$.
  \source{petersen:page-0416}
\end{lemma}
\begin{proof}
  \uses{lem:pet-ch11-gh-close-cov-cap-comparison}
  Fix an admissible metric on $X\cup Y$ realizing $d_H(X,Y)<\delta$. A
  $\Cov_Y(\varepsilon)$-cover of $Y$ by $\varepsilon$-balls pulls back, via
  the triangle inequality in the admissible metric, to a cover of $X$ by
  $(\varepsilon+2\delta)$-balls, giving the first inequality; a disjoint
  family of $(\varepsilon+2\delta)$-balls in $Y$ likewise gives, by the same
  triangle-inequality argument, a disjoint family of $\varepsilon$-balls in
  $X$, giving the second.
\end{proof}

\begin{proposition}[Prop.\ 11.1.10 (Gromov, 1980) — precompactness criterion]
  \label{prop:pet-ch11-gromov-precompactness-criterion}
  \lean{PetersenLib.gromovPrecompactnessCriterion}
  \uses{lem:pet-ch11-cov-le-cap,lem:pet-ch11-gh-close-cov-cap-comparison}
  For $\mathcal{C}\subset(\mathcal{M},\dGH)$ with all diameters $\le D<\infty$,
  the following are equivalent: (1) $\mathcal{C}$ is precompact; (2) there is
  $N_1:(0,\infty)\to(0,\infty)$ with $\Cap_X(\varepsilon)\le N_1(\varepsilon)$
  for all $X\in\mathcal{C}$; (3) there is $N_2:(0,\infty)\to(0,\infty)$ with
  $\Cov_X(\varepsilon)\le N_2(\varepsilon)$ for all $X\in\mathcal{C}$.
  \source{petersen:page-0416,petersen:page-0417}
\end{proposition}
\begin{proof}
  \uses{prop:pet-ch11-gromov-precompactness-criterion}
  (1)$\Rightarrow$(2): precompactness gives $X_1,\dots,X_k\in\mathcal{C}$ with
  $\dGH(X,X_i)<\varepsilon/4$ for some $i$, for every $X\in\mathcal{C}$; by
  \cref{lem:pet-ch11-gh-close-cov-cap-comparison},
  $\Cap_X(\varepsilon)\le\Cap_{X_i}(\varepsilon/2)\le\max_i\Cap_{X_i}(\varepsilon/2)$.
  (2)$\Rightarrow$(3): take $N_2=N_1$ and use
  \cref{lem:pet-ch11-cov-le-cap}. (3)$\Rightarrow$(1): it suffices to show
  $\mathcal{C}$ is totally bounded. Since
  $\Cov_X(\varepsilon/2)\le N(\varepsilon/2)$, every $X\in\mathcal{C}$ is
  $\varepsilon/2$-close to a finite subset of at most $N(\varepsilon/2)$
  points, giving a finite metric space with distance matrix
  $(d_{ij})\in[0,D]^{N(\varepsilon/2)\times N(\varepsilon/2)}$; a finite
  $\varepsilon/2$-net in this compact space of matrices produces finitely
  many finite metric spaces such that every such matrix is within
  $\varepsilon/2$ of one of them (entrywise), and hence every $X\in\mathcal{C}$
  is within $\varepsilon$ of one of the finitely many chosen spaces.
\end{proof}

\begin{corollary}[Cor.\ 11.1.11 — pointed precompactness]
  \label{cor:pet-ch11-pointed-precompactness-criterion}
  \lean{PetersenLib.pointedPrecompactnessCriterion}
  \uses{prop:pet-ch11-gromov-precompactness-criterion,def:pet-ch11-pointed-gh-topology}
  $\mathcal{C}\subset\mathcal{M}_*$ is precompact iff for each $R>0$ the
  collection $\{\overline{B}(x,R)\mid(X,x)\in\mathcal{C}\}\subset(\mathcal{M},\dGH)$
  is precompact.
  \source{petersen:page-0417}
\end{corollary}

\begin{definition}[metric doubling condition]
  \label{def:pet-ch11-metric-doubling-condition}
  \lean{PetersenLib.MetricDoublingCondition}
  $X$ satisfies the \emph{metric doubling condition} with constant $C$ if
  every ball $B(p,R)$ can be covered by at most $C$ balls of radius $R/2$.
  \source{petersen:page-0417}
\end{definition}

\begin{proposition}[Prop.\ 11.1.12 — doubling implies precompact]
  \label{prop:pet-ch11-doubling-implies-precompact}
  \lean{PetersenLib.doublingCondition_implies_precompact}
  \uses{def:pet-ch11-metric-doubling-condition,prop:pet-ch11-gromov-precompactness-criterion}
  If every $X\in\mathcal{C}\subset(\mathcal{M},\dGH)$ satisfies the metric
  doubling condition with constant $C<\infty$ and $\diam X\le D<\infty$, then
  $\mathcal{C}$ is precompact.
  \source{petersen:page-0417}
\end{proposition}
\begin{proof}
  \uses{prop:pet-ch11-doubling-implies-precompact}
  Iterating the doubling condition, every $X\in\mathcal{C}$ is covered by at
  most $C^N$ balls of radius $2^{-N}D$, hence by at most $C^N$ balls of any
  radius $\varepsilon\in[2^{-N}D,2^{-N+1}D]$. This gives the required bound
  $\Cov_X(\varepsilon)\le N_2(\varepsilon)$, and
  \cref{prop:pet-ch11-gromov-precompactness-criterion} applies.
\end{proof}

\begin{corollary}[Cor.\ 11.1.13 — Ricci/diameter precompactness]
  \label{cor:pet-ch11-ricci-diam-precompactness}
  \lean{PetersenLib.riccciBounded_precompact}
  \uses{prop:pet-ch11-doubling-implies-precompact,cor:pet-ch11-pointed-precompactness-criterion}
  For $n\ge2$, $k\in\mathbb{R}$, $D>0$: (1) the class of closed Riemannian
  $n$-manifolds with $\Ric\ge(n-1)k$, $\diam\le D$ is precompact; (2) the
  class of pointed complete Riemannian $n$-manifolds with $\Ric\ge(n-1)k$ is
  precompact.
  \source{petersen:page-0417}
\end{corollary}
\begin{proof}
  \uses{cor:pet-ch11-ricci-diam-precompactness}
  It suffices to prove (2). Fix $R>0$; disjoint balls
  $B(x_1,\varepsilon),\dots,B(x_N,\varepsilon)\subset\overline{B}(x,R)$ have,
  taking $B(x_i,\varepsilon)$ of minimal volume,
  \[
    N\le\frac{\vol B(x,R)}{\vol B(x_i,\varepsilon)}\le\frac{v(n,k,2R)}{v(n,k,\varepsilon)}
  \]
  by relative volume comparison, giving a bound on the number of disjoint
  $\varepsilon$-balls in $\overline{B}(x,R)$ independent of the manifold, so
  \cref{cor:pet-ch11-pointed-precompactness-criterion} applies via
  \cref{prop:pet-ch11-doubling-implies-precompact}-style capacity control.
\end{proof}

\section{Hölder Spaces and Schauder Estimates}

\subsection{Hölder Spaces}

\begin{definition}[$C^0$ and $C^m$ norms]
  \label{def:pet-ch11-cm-norm}
  \lean{PetersenLib.cZeroNorm,PetersenLib.cmNorm}
  For a bounded domain $\Omega\subset\mathbb{R}^n$, the bounded continuous
  functions $u:\Omega\to\mathbb{R}^k$ form $C^0(\Omega,\mathbb{R}^k)$ with sup-norm
  $\|u\|_{C^0}=\sup_{x\in\Omega}|u(x)|$, a Banach space. Using multi-index
  notation $\partial^Iu=\partial_1^{i_1}\cdots\partial_n^{i_n}u$ for
  $I=(i_1,\dots,i_n)$, $|I|=i_1+\dots+i_n$, define $C^m(\Omega,\mathbb{R}^k)$ as
  the functions with $m$ continuous partial derivatives, with norm
  \[
    \|u\|_{C^m}=\|u\|_{C^0}+\sum_{1\le|I|\le m}\|\partial^Iu\|_{C^0}.
  \]
  This makes $C^m(\Omega,\mathbb{R}^k)$ a Banach space, but the inclusions
  $C^m\subset C^{m-1}$ are not closed subspaces (e.g.\ $f(x)=|x|$ lies in the
  $C^0([-1,1],\mathbb{R})$-closure of $C^1([-1,1],\mathbb{R})$).
  \source{petersen:page-0418,petersen:page-0419}
\end{definition}

\begin{definition}[Hölder pseudo-norm]
  \label{def:pet-ch11-holder-pseudonorm}
  \lean{PetersenLib.holderPseudonorm}
  \uses{def:pet-ch11-cm-norm}
  For $\alpha\in(0,1]$ and $u:\Omega\to\mathbb{R}^k$, the \emph{$C^\alpha$
  pseudo-norm} is
  \[
    \|u\|_\alpha=\sup_{x,y\in\Omega}\frac{|u(x)-u(y)|}{|x-y|^\alpha}.
  \]
  When $\alpha=1$ this is the best Lipschitz constant for $u$.
  \source{petersen:page-0419}
\end{definition}

\begin{definition}[Hölder space $C^{m,\alpha}$]
  \label{def:pet-ch11-holder-space}
  \lean{PetersenLib.holderSpace}
  \uses{def:pet-ch11-cm-norm,def:pet-ch11-holder-pseudonorm}
  The \emph{Hölder space} $C^{m,\alpha}(\Omega,\mathbb{R}^k)$ consists of the
  $u\in C^m(\Omega,\mathbb{R}^k)$ whose $m$th-order partial derivatives all have
  finite $C^\alpha$ pseudo-norm, with norm
  \[
    \|u\|_{C^{m,\alpha}}=\|u\|_{C^m}+\sum_{|I|=m}\|\partial^Iu\|_\alpha.
  \]
  If the domain needs to be specified, one writes $\|u\|_{C^{m,\alpha},\Omega}$.
  \source{petersen:page-0419}
\end{definition}

\begin{lemma}[Lemma 11.2.1 — Hölder spaces are Banach, compactly nested]
  \label{lem:pet-ch11-holder-space-banach-compact-inclusion}
  \lean{PetersenLib.holderSpace_banach,PetersenLib.holderSpace_inclusion_compact}
  \uses{def:pet-ch11-holder-space}
  $C^{m,\alpha}(\Omega,\mathbb{R}^k)$ is a Banach space, and for $\beta<\alpha$ the
  inclusion $C^{m,\alpha}(\Omega,\mathbb{R}^k)\subset C^{m,\beta}(\Omega,\mathbb{R}^k)$ is
  always compact (maps closed bounded sets to compact sets).
  \source{petersen:page-0419,petersen:page-0420}
\end{lemma}
\begin{proof}
  \uses{lem:pet-ch11-holder-space-banach-compact-inclusion}
  It suffices to treat $m=0$; the general case follows immediately.

  \emph{Banach.} Let $\{u_i\}$ be Cauchy in $C^\alpha(\Omega,\mathbb{R}^k)$; it is
  Cauchy in $C^0$, so $u_i\to u\in C^0$ uniformly. For fixed $x\ne y$,
  $\frac{|u_i(x)-u_i(y)|}{|x-y|^\alpha}\to\frac{|u(x)-u(y)|}{|x-y|^\alpha}$;
  the left side is uniformly bounded, so the limit is too, giving
  $u\in C^\alpha$. Given $\varepsilon>0$ and $N$ with, for $i,j\ge N$,
  $x\ne y$,
  \[
    \frac{|(u_i(x)-u_j(x))-(u_i(y)-u_j(y))|}{|x-y|^\alpha}\le\varepsilon,
  \]
  letting $j\to\infty$ gives the same bound with $u_j$ replaced by $u$, so
  $u_i\to u$ in $C^\alpha$.

  \emph{Compact inclusion.} A bounded sequence in $C^\alpha(\Omega,\mathbb{R}^k)$
  is equicontinuous, so by Arzelà--Ascoli the inclusion
  $C^\alpha(\Omega,\mathbb{R}^k)\subset C^0(\Omega,\mathbb{R}^k)$ is compact. Using
  \[
    \frac{|u(x)-u(y)|}{|x-y|^\beta}
    =\Big(\frac{|u(x)-u(y)|}{|x-y|^\alpha}\Big)^{\beta/\alpha}\cdot|u(x)-u(y)|^{1-\beta/\alpha}
  \]
  gives $\|u\|_\beta\le(\|u\|_\alpha)^{\beta/\alpha}\cdot(2\|u\|_{C^0})^{1-\beta/\alpha}$,
  so a sequence convergent in $C^0$ and bounded in $C^\alpha$ also converges
  in $C^\beta$ for any $\beta<\alpha\le1$.
\end{proof}

\subsection{Elliptic Estimates}

\begin{definition}[elliptic operator]
  \label{def:pet-ch11-elliptic-operator}
  \lean{PetersenLib.EllipticOperator}
  \uses{def:pet-ch11-holder-pseudonorm}
  An operator $Lu=a^{ij}\partial_i\partial_ju+b^i\partial_iu$ ($a^{ij}=a^{ji}$)
  is \emph{elliptic} when $(a^{ij})$ is positive definite. Throughout, all
  eigenvalues of $(a^{ij})$ lie in $[\lambda,\lambda^{-1}]$, $\lambda>0$, and
  $\|a^{ij}\|_\alpha\le\lambda^{-1}$, $\|b^i\|_\alpha\le\lambda^{-1}$.
  \source{petersen:page-0420,petersen:page-0421}
\end{definition}

\begin{theorem}[Thm.\ 11.2.2 — Schauder estimates]
  \label{thm:pet-ch11-schauder-estimates}
  \lean{PetersenLib.schauderEstimates}
  \uses{def:pet-ch11-elliptic-operator,def:pet-ch11-holder-space}
  Let $\Omega\subset\mathbb{R}^n$ have diameter $\le D$ and $K\subset\Omega$ with
  $d(K,\partial\Omega)\ge\delta$. For $\alpha\in(0,1)$ there is
  $C=C(n,\alpha,\lambda,\delta,D)$ with
  \[
    \|u\|_{C^{2,\alpha},K}\le C\big(\|Lu\|_{C^0,\Omega}+\|u\|_{C^0,\Omega}\big),
    \qquad
    \|u\|_{C^{1,\alpha},K}\le C\big(\|Lu\|_{C^0,\Omega}+\|u\|_{C^0,\Omega}\big).
  \]
  If in addition $\Omega$ has smooth boundary and $u=\varphi$ on
  $\partial\Omega$, there is $C=C(n,\alpha,\lambda,D)$ with
  \[
    \|u\|_{C^{2,\alpha},\Omega}\le C\big(\|Lu\|_{C^0,\Omega}+\|\varphi\|_{C^{2,\alpha},\partial\Omega}\big).
  \]
  \source{petersen:page-0421}
\end{theorem}

\begin{remark}[divergence form and the Laplace--Beltrami operator]
  \label{rem:pet-ch11-divergence-form-laplace-beltrami}
  \lean{PetersenLib.divergenceFormOperator,PetersenLib.laplaceBeltrami}
  \uses{def:pet-ch11-elliptic-operator}
  One proves \cref{thm:pet-ch11-schauder-estimates} first for $Lu=\Delta u=\delta^{ij}\partial_i\partial_ju$,
  then, via a linear change of coordinates, for constant-coefficient
  operators $Lu=a^{ij}\partial_i\partial_ju$; Schauder's trick is that the
  hypotheses on $a^{ij}$ make it locally almost constant, and a partition of
  unity finishes the argument. The first-order term causes no extra trouble,
  especially when $L$ is in \emph{divergence form}
  $Lu=a^{ij}\partial_i\partial_ju+b^i\partial_iu=\partial_i(a^{ij}\partial_ju)$,
  which integrates by parts as
  $\int_\Omega\partial_i(a^{ij}\partial_ju)\,h=-\int_\Omega a^{ij}\partial_iu\,\partial_ih$
  for $h=0$ on $\partial\Omega$. The Laplacian of a Riemannian metric $g$ is of
  this form in local coordinates:
  \[
    Lu=\Delta_gu=\frac1{\sqrt{\det g_{ij}}}\partial_i\big(\sqrt{\det g_{ij}}\,g^{ij}\partial_ju\big),
  \]
  so $\int v\,Lu\,\vol=\int v\,\partial_i\big(\sqrt{\det g_{ij}}\,g^{ij}\partial_ju\big)$.
  \source{petersen:page-0421,petersen:page-0422}
\end{remark}

\begin{corollary}[Cor.\ 11.2.3 — higher-order Schauder estimates]
  \label{cor:pet-ch11-higher-schauder-estimates}
  \lean{PetersenLib.schauderEstimates_higherOrder}
  \uses{thm:pet-ch11-schauder-estimates}
  If also $\|a^{ij}\|_{C^{m,\alpha}},\|b^i\|_{C^{m,\alpha}}\le\lambda^{-1}$,
  there is $C=C(n,m,\alpha,\lambda,\delta,D)$ with
  \[
    \|u\|_{C^{m+2,\alpha},K}\le C\big(\|Lu\|_{C^{m,\alpha},\Omega}+\|u\|_{C^0,\Omega}\big),
  \]
  and on a domain with smooth boundary,
  $\|u\|_{C^{m+2,\alpha},\Omega}\le C\big(\|Lu\|_{C^{m,\alpha},\Omega}+\|\varphi\|_{C^{m+2,\alpha},\partial\Omega}\big)$.
  \source{petersen:page-0422}
\end{corollary}

\begin{theorem}[Thm.\ 11.2.4 — Dirichlet problem]
  \label{thm:pet-ch11-dirichlet-problem-unique-solution}
  \lean{PetersenLib.dirichletProblem_existsUnique}
  \uses{thm:pet-ch11-schauder-estimates}
  If $\Omega\subset\mathbb{R}^n$ is a bounded domain with smooth boundary, the
  Dirichlet problem $Lu=f$, $u|_{\partial\Omega}=\varphi$ always has a unique
  solution $u\in C^{2,\alpha}(\Omega)$ when $f\in C^\alpha(\Omega)$,
  $\varphi\in C^{2,\alpha}(\partial\Omega)$.
  \source{petersen:page-0422}
\end{theorem}
\begin{proof}
  \uses{thm:pet-ch11-dirichlet-problem-unique-solution}
  Uniqueness is an immediate consequence of the maximum principle: if $Lu_1=Lu_2=f$
  with $u_1=u_2=\varphi$ on $\partial\Omega$, then $w=u_1-u_2$ solves
  $Lw=0$ with $w=0$ on $\partial\Omega$, and the maximum principle for the
  elliptic operator $L$ forces $w\equiv0$ on $\Omega$, i.e.\ $u_1=u_2$.
\end{proof}

\subsection{Harmonic Coordinates}

\begin{lemma}[Lemma 11.2.5 — existence of harmonic coordinates]
  \label{lem:pet-ch11-harmonic-coordinates-exist}
  \lean{PetersenLib.harmonicCoordinates_exist}
  \uses{thm:pet-ch11-schauder-estimates,rem:pet-ch11-divergence-form-laplace-beltrami}
  If $(M,g)$ is an $n$-dimensional Riemannian manifold and $p\in M$, there is
  a neighborhood $U\ni p$ carrying a \emph{harmonic coordinate system}
  $x=(x^1,\dots,x^n):U\to\mathbb{R}^n$, i.e.\ each $x^i$ is harmonic for
  $\Delta_g$.
  \source{petersen:page-0422,petersen:page-0423}
\end{lemma}
\begin{proof}
  \uses{lem:pet-ch11-harmonic-coordinates-exist}
  Choose coordinates $y=(y^1,\dots,y^n)$ near $p$ with $y(p)=0$, identifying
  $M$ locally with an open subset of $\mathbb{R}^n$ and $p=0$; write
  $g_{ij}=g(\partial_i,\partial_j)$ in these Cartesian coordinates. Seek
  $y\mapsto x$ with $\Delta x^k=0$ for all $k$. Fix a small ball
  $B(0,\varepsilon)$ and solve the Dirichlet problem $\Delta x^k=0$,
  $x^k=y^k$ on $\partial B(0,\varepsilon)$
  (\cref{thm:pet-ch11-dirichlet-problem-unique-solution}), giving $n$
  harmonic functions. By \cref{thm:pet-ch11-schauder-estimates},
  \[
    \|x-y\|_{C^{2,\alpha},B(0,\varepsilon)}
    \le C\big(\|\Delta(x-y)\|_{C^0,B(0,\varepsilon)}+\|x-y\|_{C^{2,\alpha},\partial B(0,\varepsilon)}\big)
    =C\|\Delta y\|_{C^\alpha,B(0,\varepsilon)},
  \]
  using $\Delta x=0$ and $x=y$ on $\partial B(0,\varepsilon)$. Choosing $y$ to
  be exponential (normal) coordinates gives $\partial_kg_{ij}(0)=0$, so the
  formula for $\Delta_g$ in \cref{rem:pet-ch11-divergence-form-laplace-beltrami}
  gives $\Delta y(0)=0$, and the constant $C$ depends only on an upper bound
  for $\diam B(0,\varepsilon)$ besides $\alpha,n,\lambda$; hence
  $\|\Delta y\|_{C^\alpha,B(0,\varepsilon)}\to0$ and
  $\|x-y\|_{C^{2,\alpha},B(0,\varepsilon)}\to0$ as $\varepsilon\to0$. Thus for
  small $\varepsilon$, $Dx$ is close to $Dy$, which is invertible since $y$ is
  a coordinate system; hence $x$ is a coordinate system too.
\end{proof}

\begin{lemma}[Lemma 11.2.6 — Laplacian and Ricci curvature in harmonic coordinates]
  \label{lem:pet-ch11-laplacian-ricci-harmonic-coords}
  \lean{PetersenLib.laplacian_eq_gijPartial_harmonic,PetersenLib.ricci_harmonicCoords_formula}
  \uses{lem:pet-ch11-harmonic-coordinates-exist}
  On a harmonic coordinate system $x:U\to\mathbb{R}^n$ for $(M,g)$,
  \begin{equation}
    \Delta u=\frac1{\sqrt{\det g_{ij}}}\partial_i\big(\sqrt{\det g_{ij}}\,g^{ij}\partial_ju\big)=g^{ij}\partial_i\partial_ju,
  \end{equation}
  \begin{equation}
    \tfrac12\Delta g_{ij}+Q(g,\partial g)=-\Ric_{ij}=-\Ric(\partial_i,\partial_j),
  \end{equation}
  where $Q$ is a universal rational expression, polynomial in $g$ and
  quadratic in $\partial g$ in the numerator, with denominator a power of
  $\sqrt{\det g_{ij}}$.
  \source{petersen:page-0423,petersen:page-0424,petersen:page-0425}
\end{lemma}
\begin{proof}
  \uses{lem:pet-ch11-laplacian-ricci-harmonic-coords}
  \emph{(1)} Since $x^k$ is harmonic,
  \[
    0=\Delta x^k=g^{ij}\partial_i\partial_jx^k
    +\frac1{\sqrt{\det g_{ij}}}\partial_i\big(\sqrt{\det g_{ij}}\,g^{ij}\big)\partial_jx^k
    =\frac1{\sqrt{\det g_{ij}}}\partial_i\big(\sqrt{\det g_{ij}}\,g^{ik}\big),
  \]
  using $\partial_i\partial_jx^k=0$ and $\partial_jx^k=\delta_j^k$ in these
  coordinates. Hence in
  $\Delta u=g^{ij}\partial_i\partial_ju+\frac1{\sqrt{\det g_{ij}}}\partial_i(\sqrt{\det g_{ij}}\,g^{ij})\partial_ju$
  the second term vanishes, giving $\Delta u=g^{ij}\partial_i\partial_ju$.

  \emph{(2)} For harmonic $u$ the Bochner formula gives
  $\Delta(\tfrac12|\nabla u|^2)=|\Hess u|^2+\Ric(\nabla u,\nabla u)$, where
  $|\Hess u|^2$ depends only on the metric and its first derivatives.
  Applying this to $u=x^k$ and using $g(\nabla x^i,\nabla x^j)=g^{ij}$
  (since $\nabla x^k=g^{ik}\partial_i$) gives
  \[
    \tfrac12\Delta g^{ij}-g(\Hess x^i,\Hess x^j)=\Ric(\nabla x^i,\nabla x^j)
  \]
  after polarizing the quadratic identity
  $\tfrac12\Delta g(\nabla x^k,\nabla x^k)-|\Hess x^k|^2=\Ric(\nabla x^k,\nabla x^k)$.
  Differentiating the matrix identity $[g_{ik}][g^{kj}]=[\delta_i^j]$ twice
  gives
  \[
    0=[\Delta g_{ik}][g^{kj}]+2[\nabla g_{ik}]\cdot[\nabla g^{kj}]+[g_{ik}][\Delta g^{kj}],
  \]
  and substituting the formula for $\Delta g^{kj}$ from the previous display
  and solving for $[\Delta g_{ij}]$ yields
  \[
    \tfrac12\Delta g_{ij}+Q_{ij}(g,\partial g)=-\Ric_{ij},\qquad
    Q_{ij}=-2\sum_{k,l}g(\nabla g_{ik},\nabla g^{kl})g_{lj}-2\sum_{k,l}g_{ik}\,g(\Hess x^k,\Hess x^l)\,g_{lj},
  \]
  which is polynomial in $g$, quadratic in $\partial g$, and rational with
  denominator a power of $\sqrt{\det g_{ij}}$ (via $g^{ij}$), as claimed.
\end{proof}

\begin{remark}[Einstein metrics in harmonic coordinates are smooth]
  \label{rem:pet-ch11-einstein-harmonic-regularity-bootstrap}
  \lean{PetersenLib.einsteinHarmonicWeakSolution_smooth}
  \uses{lem:pet-ch11-laplacian-ricci-harmonic-coords,cor:pet-ch11-higher-schauder-estimates}
  For an Einstein metric $\Ric_{ij}=(n-1)kg_{ij}$,
  \cref{lem:pet-ch11-laplacian-ricci-harmonic-coords} reads
  $\tfrac12\Delta g_{ij}=-(n-1)kg_{ij}-Q(g,\partial g)$, which makes sense
  once $g_{ij}\in C^1$ and, after multiplying by a test function and
  integrating by parts, can be interpreted weakly. If $g_{ij}\in C^{1,\alpha}$
  the right side lies in some $C^\beta$, so \cref{cor:pet-ch11-higher-schauder-estimates}
  gives $g_{ij}\in C^{2,\beta}$; bootstrapping shows $g_{ij}\in C^\infty$
  (in fact real-analytic). Hence any metric that is, in harmonic coordinates,
  a weak solution of the Einstein equation is automatically smooth.
  \source{petersen:page-0425,petersen:page-0426}
\end{remark}

\section{Norms and Convergence of Manifolds}

\subsection{Norms of Riemannian Manifolds}

\begin{definition}[$C^{m,\alpha}$ norm at scale $r$]
  \label{def:pet-ch11-cmalpha-norm-pointed}
  \lean{PetersenLib.pointedNorm}
  \uses{def:pet-ch11-holder-space}
  For a pointed Riemannian $n$-manifold $(M,g,p)$, say
  $\|(M,g,p)\|_{C^{m,\alpha},r}\le Q$ if there is a $C^{m+1,\alpha}$ chart
  $\varphi:(B(0,r),0)\subset\mathbb{R}^n\to(U,p)\subset M$ such that
  \begin{align}
    \text{(n1)}\quad&|D\varphi|\le e^Q\text{ on }B(0,r),\ |D\varphi^{-1}|\le e^Q\text{ on }U,
    \text{ equivalently }e^{-2Q}\delta_{kl}v^kv^l\le g_{kl}v^kv^l\le e^{2Q}\delta_{kl}v^kv^l\ \forall v\in\mathbb{R}^n,\\
    \text{(n2)}\quad&r^{|I|+\alpha}\|\partial^Ig_{ij}\|_\alpha\le Q\text{ for all multi-indices }I,\ 0\le|I|\le m.
  \end{align}
  \source{petersen:page-0426,petersen:page-0427}
\end{definition}

\begin{definition}[global $C^{m,\alpha}$ norm of a manifold]
  \label{def:pet-ch11-cmalpha-norm-global}
  \lean{PetersenLib.globalNorm}
  \uses{def:pet-ch11-cmalpha-norm-pointed}
  $\|(M,g)\|_{C^{m,\alpha},r}=\sup_{p\in M}\|(M,g,p)\|_{C^{m,\alpha},r}$. The
  charts are regarded as maps from the fixed model $B(0,r)$ into $M$, so that
  the norm measures how far the manifold differs from a Euclidean ball: (n1)
  bounds the metric coefficients from above and below in the chosen
  coordinates (giving uniform ellipticity for the Laplacian), and (n2)
  bounds the derivatives of the metric. The definition makes sense for
  merely $C^{m,\alpha}$ metrics, without assuming more smoothness; some
  constructions (e.g.\ exponential maps) may then be ill-defined if $m\le1$.
  The pointed norm is always finite, but the global norm need not be finite
  on any scale when $M$ is noncompact.
  \source{petersen:page-0427}
\end{definition}

\begin{example}[Example 11.3.1 — flat manifolds have norm zero]
  \label{ex:pet-ch11-flat-manifold-norm-zero}
  \lean{PetersenLib.flatManifold_norm_zero}
  \uses{def:pet-ch11-cmalpha-norm-global}
  If $(M,g)$ is a complete flat manifold then $\|(M,g)\|_{C^{m,\alpha},r}=0$
  for all $r\le\inj(M,g)$; in particular
  $\|(\mathbb{R}^n,g_{\mathrm{eu}})\|_{C^{m,\alpha},r}=0$ for all $r$.
  \source{petersen:page-0427}
\end{example}

\subsection{Convergence of Riemannian Manifolds}

\begin{definition}[$C^{m,\alpha}$ convergence of tensors on a compact subset]
  \label{def:pet-ch11-tensor-convergence-compact-subset}
  \lean{PetersenLib.tensorConverges_onCompactSubset}
  Fix a closed manifold $M$, or a precompact subset $A\subset M$. A sequence
  of functions on $A$ \emph{converges in $C^{m,\alpha}$} if it converges in
  the charts of some fixed finite covering of coordinate patches that are
  uniformly bi-Lipschitz; this is independent of the covering. A sequence of
  tensors converges in $C^{m,\alpha}$ if its components converge in these
  patches, giving a notion of convergence of Riemannian metrics on compact
  subsets of a fixed manifold.
  \source{petersen:page-0427}
\end{definition}

\begin{definition}[pointed $C^{m,\alpha}$ convergence of manifolds]
  \label{def:pet-ch11-pointed-cmalpha-convergence}
  \lean{PetersenLib.PointedCmAlphaConvergence}
  \uses{def:pet-ch11-tensor-convergence-compact-subset,def:pet-ch11-cmalpha-norm-pointed}
  A sequence of pointed complete Riemannian manifolds converges in the
  \emph{pointed $C^{m,\alpha}$ topology}, $(M_i,g_i,p_i)\to(M,g,p)$, if for
  every $R>0$ there is a domain $\Omega\supset B(p,R)\subset M$ and
  embeddings $F_i:\Omega\to M_i$ for large $i$ with $F_i(p)=p_i$,
  $F_i(\Omega)\supset B(p_i,R)$, and $F_i^*g_i\to g$ on $\Omega$ in the
  $C^{m,\alpha}$ topology.
  \source{petersen:page-0427,petersen:page-0428}
\end{definition}

\begin{remark}[pointed $C^{m,\alpha}$ convergence implies GH convergence]
  \label{rem:pet-ch11-cmalpha-convergence-implies-gh}
  \lean{PetersenLib.cmAlphaConvergence_implies_ghConvergence}
  \uses{def:pet-ch11-pointed-cmalpha-convergence,def:pet-ch11-pointed-gh-topology}
  Pointed $C^{m,\alpha}$ convergence implies pointed Gromov--Hausdorff
  convergence. When all manifolds are closed with uniformly bounded
  diameter, the $F_i$ are diffeomorphisms, so one may speak of unpointed
  convergence, which forces all manifolds in the tail of the sequence to be
  diffeomorphic; classes of closed manifolds precompact in some $C^{m,\alpha}$
  topology contain at most finitely many diffeomorphism types.
  \source{petersen:page-0428}
\end{remark}

\begin{remark}[pullback convergence versus classical convergence]
  \label{rem:pet-ch11-cmalpha-convergence-pullback-warning}
  \lean{PetersenLib.pullbackConvergence_not_classicalConvergence}
  \uses{def:pet-ch11-pointed-cmalpha-convergence}
  Metrics $g_i=F_i^*g$ for diffeomorphisms $F_i$ of a fixed $M$ can converge
  in the pointed $C^{m,\alpha}$ sense of
  \cref{def:pet-ch11-pointed-cmalpha-convergence} without converging in any
  fixed coordinate system: e.g., on $M=S^2$ with the standard metric $g$, let
  $F_i$ be diffeomorphisms from the flow of a vector field vanishing only at
  the poles and pointing toward the south pole. Away from the poles,
  $F_i^*g\to0$ in fixed coordinates as the flow collapses the sphere toward
  the south pole, so $(S^2,F_i^*g)$ does not converge classically; but
  $(S^2,F_i^*g)$ is isometric to $(S^2,g)$ for every $i$, so as pointed
  Riemannian manifolds up to isometry the sequence is constant.
  \source{petersen:page-0428}
\end{remark}

\subsection{Properties of the Norm}

\begin{proposition}[Prop.\ 11.3.2 — elementary properties of the norm]
  \label{prop:pet-ch11-norm-properties}
  \lean{PetersenLib.norm_scaling,PetersenLib.norm_monotone_continuous,PetersenLib.norm_continuous_underConvergence,PetersenLib.norm_ballContainment,PetersenLib.norm_chart_attained,PetersenLib.norm_sup_attained}
  \uses{def:pet-ch11-cmalpha-norm-pointed,def:pet-ch11-cmalpha-norm-global,def:pet-ch11-pointed-cmalpha-convergence}
  Given $(M,g,p)$, $m\ge0$, $\alpha\in(0,1]$:
  \begin{enumerate}
  \item[(1)] $\|(M,g,p)\|_{C^{m,\alpha},r}=\|(M,\lambda^2g,p)\|_{C^{m,\alpha},\lambda r}$ for all $\lambda>0$.
  \item[(2)] $r\mapsto\|(M,g,p)\|_{C^{m,\alpha},r}$ is increasing, continuous, and $\to0$ as $r\to0$.
  \item[(3)] If $(M_i,g_i,p_i)\to(M,g,p)$ in $C^{m,\alpha}$, then
    $\|(M_i,g_i,p_i)\|_{C^{m,\alpha},r}\to\|(M,g,p)\|_{C^{m,\alpha},r}$ for all $r>0$;
    if in addition all manifolds have uniformly bounded diameter,
    $\|(M_i,g_i)\|_{C^{m,\alpha},r}\to\|(M,g)\|_{C^{m,\alpha},r}$ for all $r>0$.
  \item[(4)] If $\|(M,g,p)\|_{C^{m,\alpha},r}<Q$, then for $x_1,x_2\in B(0,r)$,
    $e^{-Q}\min\{|x_1-x_2|,2r-|x_1|-|x_2|\}\le|\varphi(x_1)\varphi(x_2)|\le e^Q|x_1-x_2|$.
  \item[(5)] The infimum $Q=\|(M,g,p)\|_{C^{m,\alpha},r}$ is attained by some
    chart satisfying (n1), (n2) with this $Q$.
  \item[(6)] When $M$ is compact, the supremum $\|(M,g)\|_{C^{m,\alpha},r}=\sup_p\|(M,g,p)\|_{C^{m,\alpha},r}$ is a maximum.
  \end{enumerate}
  \source{petersen:page-0428,petersen:page-0429,petersen:page-0430,petersen:page-0431}
\end{proposition}
\begin{proof}
  \uses{prop:pet-ch11-norm-properties}
  \emph{(1)} Given a chart $\varphi:B(0,r)\to M$ for $g$, set
  $\varphi^*(x)=\varphi(\lambda^{-1}x):B(0,\lambda r)\to M$; scaling the metric
  by $\lambda^2$ leaves (n1), (n2) valid with the same $Q$.

  \emph{(2)} Restricting a chart $\varphi:B(0,r)\to M$ to a smaller ball shows
  $N(r):=\|(M,g,p)\|_{C^{m,\alpha},r}$ is increasing. Considering
  $\varphi^*(x)=\varphi(\lambda^{-1}x):B(0,\lambda r)\to M$ without rescaling
  $g$, one gets from $N(r)<Q$ that
  $N(\lambda r)\le\max\{Q\pm|\log\lambda|,\,Q\lambda^2\}$. With $\lambda=r_i/r$,
  $r_i\to r$, this gives $\limsup N(r_i)\le N(r)$; applying the same estimate
  with $r,r_i$ interchanged gives $N(r)\le\liminf N(r_i)$, so $N$ is
  continuous. For $N(r)\to0$ as $r\to0$: any chart around $p$ can, after a
  linear change, be arranged so $g_{ij}(p)=\delta_{ij}$, hence
  $|D\varphi|_p|=|D\varphi^{-1}|_p|=1$; using this chart on sufficiently
  small balls gives $N(r)\to0$.

  \emph{(3)} Fix $r>0$, $Q>\|(M,g,p)\|_{C^{m,\alpha},r}$, and a domain
  $\Omega\supset B(p,e^Qr)$ with embeddings $F_i:\Omega\to M_i$,
  $F_i^*g_i\to g$ in $C^{m,\alpha}$, $F_i(p)=p_i$, for large $i$. Given a
  chart $\varphi:B(0,r)\to M$ satisfying (n1), (n2) for $Q$, the charts
  $\varphi_i=F_i\circ\varphi:B(0,r)\to M_i$ satisfy (n1), (n2) for
  $Q_i\to Q$, giving $\limsup\|(M_i,g_i,p_i)\|_{C^{m,\alpha},r}\le\|(M,g,p)\|_{C^{m,\alpha},r}$.
  Conversely, if $Q>\|(M_i,g_i,p_i)\|_{C^{m,\alpha},r}$ for large $i$, a chart
  $\varphi_i:B(0,r)\to M_i$ pulls back via $\varphi=F_i^{-1}\circ\varphi_i$ to
  give $\|(M,g,p)\|_{C^{m,\alpha},r}\le Q_i$ with $Q_i$ close to $Q$, giving
  $\liminf\|(M_i,g_i,p_i)\|_{C^{m,\alpha},r}\ge\|(M,g,p)\|_{C^{m,\alpha},r}$.
  Combined, $\|(M_i,g_i,p_i)\|_{C^{m,\alpha},r}\to\|(M,g,p)\|_{C^{m,\alpha},r}$.
  For uniformly bounded diameter, choose diffeomorphisms $F_i:M\to M_i$ for
  large $i$ with $F_i^*g_i\to g$; for each $p\in M$, $p_i=F_i(p)$, the
  pointed statement gives
  $\liminf\|(M_i,g_i)\|_{C^{m,\alpha},r}\ge\sup_p\|(M,g,p)\|_{C^{m,\alpha},r}$
  and, for $p_i\in M_i$, $p=F_i^{-1}(p_i)$,
  $\limsup\|(M_i,g_i,p_i)\|_{C^{m,\alpha},r}\le\sup_p\|(M,g)\|_{C^{m,\alpha},r}$.

  \emph{(4)} $|D\varphi|\le e^Q$ and convexity of $B(0,r)$ give the upper
  bound directly. For the lower bound, if the segment from $\varphi(x_1)$ to
  $\varphi(x_2)$ lies in $U$, then $|D\varphi^{-1}|\le e^Q$ gives
  $|x_1-x_2|\le e^Q|\varphi(x_1)\varphi(x_2)|$. Otherwise, split a
  minimizing segment $c:[0,1]\to M$ from $\varphi(x_1)$ to $\varphi(x_2)$
  into $c|_{[0,t_1)}\subset U$ and $c|_{(t_2,1]}\subset U$ with
  $c(t_i)\in\partial U$; then
  \[
    |\varphi(x_1)\varphi(x_2)|=L(c)\ge L(c|_{[0,t_1)})+L(c|_{(t_2,1]})
    \ge e^{-Q}\big(L(\varphi^{-1}\circ c|_{[0,t_1)})+L(\varphi^{-1}\circ c|_{(t_2,1]})\big)
    \ge e^{-Q}(2r-|x_1|-|x_2|),
  \]
  since $\varphi^{-1}\circ c(t)\to\partial B(0,r)$ as $t\to t_1,t_2$.

  \emph{(5)} Given charts $\varphi_i:B(0,r)\to M$ satisfying (n1), (n2) with
  $Q_i\to Q$, Arzelà--Ascoli produces a subsequence converging to a
  $C^{m+1,\alpha}$ map $\varphi:B(0,r)\to M$; part (4) shows $\varphi$ is
  injective, hence a homeomorphism onto its image, so a chart. Passing to a
  further subsequence so that the metric coefficients converge shows
  $\varphi$ satisfies (n1), (n2) for $Q$.

  \emph{(6)} Part (3) shows $p\mapsto\|(M,g,p)\|_{C^{m,\alpha},r}$ is
  continuous; compactness of $M$ then makes the supremum a maximum.
\end{proof}

\begin{corollary}[Cor.\ 11.3.3 — ball containment in the chart domain]
  \label{cor:pet-ch11-norm-bound-ball-containment}
  \lean{PetersenLib.norm_ball_subset_chartDomain}
  \uses{prop:pet-ch11-norm-properties}
  If $\|(M,g,p)\|_{C^{m,\alpha},r}\le Q$, then $B(p,e^{-Q}r)\subset U$.
  \source{petersen:page-0431}
\end{corollary}
\begin{proof}
  \uses{cor:pet-ch11-norm-bound-ball-containment}
  Let $q\in\partial U$ be closest to $p$, so $B(p,|qp|)\subset U$, and let
  $c:[0,|qp|]\to M$ be a minimizing segment from $p$ to $q$. Writing
  $c(s)=\varphi(\bar c(s))$ for $s<|qp|$, with $|\bar c(t)|\to r$ as
  $t\to|qp|$, part (4) of \cref{prop:pet-ch11-norm-properties} gives
  \[
    |qp|\ge\lim_{s\to|qp|}|\varphi(\bar c(s))\varphi(0)|
    \ge\lim_{s\to|qp|}e^{-Q}\min\{|\bar c(s)|,2r-|\bar c(s)|\}
    =e^{-Q}r.
  \]
\end{proof}

\begin{corollary}[Cor.\ 11.3.4 — norm zero implies flat]
  \label{cor:pet-ch11-norm-zero-implies-flat}
  \lean{PetersenLib.norm_zero_implies_flat}
  \uses{prop:pet-ch11-norm-properties,cor:pet-ch11-norm-bound-ball-containment}
  If $\|(M,g,p)\|_{C^{m,\alpha},r}=0$ for some $r$, then $p$ has a flat
  neighborhood.
  \source{petersen:page-0431}
\end{corollary}
\begin{proof}
  \uses{cor:pet-ch11-norm-zero-implies-flat}
  By \cref{prop:pet-ch11-norm-properties}(5) there is a $C^{m+1,\alpha}$
  chart $\varphi:B(0,r)\to U\supset B(p,e^{-Q}r)$ with $Q=0$, which is
  therefore a $C^1$ Riemannian isometry between $B(0,r)\subset\mathbb{R}^n$ and
  $U$; a $C^1$ Riemannian isometry is automatically a Riemannian isometry (a
  fact proved elsewhere in the book), so $U$ is flat.
\end{proof}

\subsection{The Harmonic Norm}

\begin{definition}[harmonic norm]
  \label{def:pet-ch11-harmonic-norm}
  \lean{PetersenLib.harmonicNorm}
  \uses{def:pet-ch11-cmalpha-norm-pointed,lem:pet-ch11-harmonic-coordinates-exist}
  The \emph{harmonic norm} $\|(M,g,p)\|^{\mathrm{har}}_{C^{m,\alpha},r}$ is
  defined exactly as in \cref{def:pet-ch11-cmalpha-norm-pointed}, except that
  $\varphi^{-1}:U\to\mathbb{R}^n$ is additionally required to be harmonic for
  $g$, i.e.\ for each $j$,
  $\frac1{\sqrt{\det[g_{ij}]}}\partial_j\big(\sqrt{\det[g_{ij}]}\,g^{ij}\big)=0$.
  \source{petersen:page-0431,petersen:page-0432}
\end{definition}

\begin{proposition}[Prop.\ 11.3.5 (Anderson, 1990) — harmonic-norm analogue]
  \label{prop:pet-ch11-harmonic-norm-properties}
  \lean{PetersenLib.harmonicNorm_properties}
  \uses{def:pet-ch11-harmonic-norm,prop:pet-ch11-norm-properties,lem:pet-ch11-harmonic-coordinates-exist}
  \cref{prop:pet-ch11-norm-properties} holds for the harmonic norm when
  $m\ge1$.
  \source{petersen:page-0432}
\end{proposition}
\begin{proof}
  \uses{prop:pet-ch11-harmonic-norm-properties}
  Only the necessary changes from the proof of
  \cref{prop:pet-ch11-norm-properties} are indicated. For part (2), that the
  norm goes to $0$ as $r\to0$: solving the Dirichlet problem as in
  \cref{lem:pet-ch11-harmonic-coordinates-exist} produces coordinates around
  every $p$ with $g_{ij}(p)=\delta_{ij}$, $\partial_kg_{ij}(p)=0$; when
  $m\ge1$, any coordinate system around $p$ can be modified to have these
  properties.

  For part (3), the inequality
  $\liminf\|(M_i,g_i,p_i)\|^{\mathrm{har}}_{C^{m,\alpha},r}\ge\|(M,g,p)\|^{\mathrm{har}}_{C^{m,\alpha},r}$:
  given $Q>\liminf\|(M_i,g_i,p_i)\|^{\mathrm{har}}_{C^{m,\alpha},r}$, take
  harmonic charts $\varphi_i:B(0,r)\to M_i$ for large $i$; after passing to a
  subsequence, $\varphi_i$ converge (via $F_i^{-1}$) to a chart
  $\varphi:B(0,r)\to M$, and since the metrics converge in $C^{m,\alpha}$ the
  Laplacians of the inverse coordinate functions converge too, so $\varphi$
  is harmonic, giving $\|(M,g,p)\|^{\mathrm{har}}_{C^{m,\alpha},r}\le Q$.

  For the reverse inequality
  $\limsup\|(M_i,g_i,p_i)\|^{\mathrm{har}}_{C^{m,\alpha},r}\le\|(M,g,p)\|^{\mathrm{har}}_{C^{m,\alpha},r}$:
  fix $Q>\|(M,g,p)\|^{\mathrm{har}}_{C^{m,\alpha},r}$ and, by continuity of
  the norm, $\varepsilon>0$ with $\|(M,g,p)\|^{\mathrm{har}}_{C^{m,\alpha},r+\varepsilon}<Q$.
  Take a chart $\varphi:B(0,r+\varepsilon)\to U\subset M$ realizing this and set
  $U_i=F_i(\varphi(B(0,r+\varepsilon/2)))\subset M_i$, a closed disc with
  smooth boundary $\partial U_i=F_i(\varphi(\partial B(0,r+\varepsilon/2)))$.
  Solve $\Delta_{g_i}\psi_i=0$ on $U_i$ with
  $\psi_i=\varphi^{-1}\circ F_i^{-1}$ on $\partial U_i$; the inverse of
  $\psi_i$, if it exists, is a coordinate chart $B(0,r)\to U_i$. Pulling
  back everything to $B(0,r+\varepsilon/2)$ via $\varphi$, comparing
  $\psi_i\circ F_i\circ\varphi$ with the identity $I$ (which agree on
  $\partial B(0,r+\varepsilon/2)$), and using that $m\ge1$ gives a $C^\alpha$
  bound on
  $\frac1{\sqrt{\det[g_i]}}\partial_k(\sqrt{\det[g_i]}\,g_i^{kl})$, one gets
  from \cref{cor:pet-ch11-higher-schauder-estimates} applied to $\Delta_{g_i}$
  on a domain with smooth boundary,
  \[
    \|I-\psi_i\circ F_i\circ\varphi\|_{C^{m+1,\alpha}}
    \le C\|\Delta_{g_i}(I-\psi_i\circ F_i\circ\varphi)\|_{C^{m-1,\alpha}}
    =C\|\Delta_{g_i}I\|_{C^{m-1,\alpha}},
  \]
  and $\|\Delta_{g_i}I\|_{C^{m-1,\alpha}}\to\|\Delta_gI\|_{C^{m-1,\alpha}}=0$
  since $F_i^*g_i\to g$ in $C^{m,\alpha}$ and $\varphi^{-1}$ is harmonic for
  $g$. Hence $\|I-\psi_i\circ F_i\circ\varphi\|_{C^{m+1,\alpha}}\to0$,
  $\psi_i$ becomes coordinates for large $i$, and these coordinates show
  $\|(M_i,g_i,p_i)\|^{\mathrm{har}}_{C^{m,\alpha},r}<Q$ for large $i$.
\end{proof}

\subsection{Compact Classes of Riemannian Manifolds}

This is the manifold analogue of the Arzelà--Ascoli lemma, essentially due to
J.\ Cheeger. It is proved in three stages: precompactness of the class in the
pointed Gromov--Hausdorff topology, the identification of the limit as a
Riemannian manifold of the same norm-bounded class, and the promotion of
Gromov--Hausdorff convergence to pointed $C^{m,\beta}$ convergence.

\begin{lemma}[Fundamental Theorem, part A — Gromov--Hausdorff precompactness]
  \label{lem:pet-ch11-fundamental-thm-precompact-gh}
  \lean{PetersenLib.fundamentalTheorem_ghPrecompact}
  \uses{def:pet-ch11-cmalpha-norm-pointed,cor:pet-ch11-ricci-diam-precompactness}
  For $Q>0$, $n\ge2$, $m\ge0$, $\alpha\in(0,1]$, $r>0$, let
  $\mathcal{M}=\mathcal{M}^{m,\alpha}(n,Q,r)$ be the class of complete,
  pointed Riemannian $n$-manifolds $(M,g,p)$ with
  $\|(M,g)\|_{C^{m,\alpha},r}\le Q$. Then $\mathcal{M}$ is precompact in the
  pointed Gromov--Hausdorff topology.
  \source{petersen:page-0434,petersen:page-0435}
\end{lemma}
\begin{proof}
  \uses{lem:pet-ch11-fundamental-thm-precompact-gh}
  Fix $M\in\mathcal{M}$, equipped with charts around every point satisfying
  (n1), (n2). Put $\delta=e^{-Q}r$. There is $N=N(n,Q)$ with $B(0,r)$
  covered by $N$ balls of radius $e^{-Q}\delta/4$; since $\varphi:B(0,r)\to U$
  is Lipschitz with constant $\le e^Q$, $U\supset B(p,\delta)$ is covered by
  $N$ balls of radius $\delta/4$. Inductively, every ball
  $B(x,\ell\delta/2)$ is covered by $\le N^\ell$ balls of radius $\delta/4$:
  for $\ell=1$ this is the previous statement; if
  $B(x,\ell\delta/2)\subset\bigcup_{i=1}^{N^\ell}B(x_i,\delta/4)$ then
  $B(x,(\ell+1)\delta/2)\subset\bigcup B(x_i,\delta)$, and covering each
  $B(x_i,\delta)$ by $\le N$ balls of radius $\delta/4$ gives $\le N^{\ell+1}$
  balls covering $B(x,(\ell+1)\delta/2)$.

  For precompactness it suffices, by
  \cref{cor:pet-ch11-pointed-precompactness-criterion} and
  \cref{prop:pet-ch11-gromov-precompactness-criterion}, to bound the number
  of disjoint $\varepsilon$-balls in $B(p,R)$ for $\varepsilon<\delta$,
  uniformly over $M\in\mathcal{M}$. Given disjoint
  $B(x_1,\varepsilon),\dots,B(x_s,\varepsilon)\subset B(p,R)$, each
  $B(x_i,\varepsilon)$ lies in a chart $\varphi:B(0,r)\to U_i$ whose preimage
  contains an $e^{-Q}\varepsilon$-ball, so
  $\vol B(x_i,\varepsilon)\ge e^{-nQ}\vol B(0,\varepsilon)$ (comparing to the
  Euclidean ball via (n1)). Taking $B(x_i,\varepsilon)$ of minimal volume,
  \[
    V(R)\ge\vol B(p,R)\ge\sum_i\vol B(x_i,\varepsilon)\ge s\cdot e^{-nQ}\vol B(0,\varepsilon),
  \]
  where $V(R)=V(R,n,Q,r)$ bounds $\vol B(p,R)$ using the ball-covering
  estimate above (each of the $\le N^\ell$ balls of radius $\delta/4$, with
  $\ell\delta/2<R\le(\ell+1)\delta/2$, has volume bounded in terms of $n,Q,r$).
  Hence $s\le C(\varepsilon)=V(R)\cdot e^{nQ}\cdot(\vol B(0,\varepsilon))^{-1}$,
  a bound depending only on $\varepsilon,R,Q,r,n$, giving the required
  capacity bound uniformly over $\mathcal{M}$.
\end{proof}

\begin{lemma}[Fundamental Theorem, part B — the limit is a norm-bounded manifold]
  \label{lem:pet-ch11-fundamental-thm-limit-is-manifold}
  \lean{PetersenLib.fundamentalTheorem_limitIsManifold}
  \uses{lem:pet-ch11-fundamental-thm-precompact-gh,prop:pet-ch11-norm-properties}
  In the setting of \cref{lem:pet-ch11-fundamental-thm-precompact-gh}, if
  $(M_i,g_i,p_i)\in\mathcal{M}$ converge to $(X,|\cdot|,p)$ in the pointed
  Gromov--Hausdorff topology, then $X$ carries the structure of a $C^{m,\alpha}$
  Riemannian manifold with $\|(X,g)\|_{C^{m,\alpha},r}\le Q$, whose induced
  distance agrees with $|\cdot|$.
  \source{petersen:page-0435}
\end{lemma}
\begin{proof}
  \uses{lem:pet-ch11-fundamental-thm-limit-is-manifold}
  For $q\in X$ choose $q_i\to q$ and charts $\varphi_i:B(0,r)\to U_i\subset M_i$
  with $q_i=\varphi_i(0)$, satisfying (n1), (n2). The $\varphi_i$ are
  uniformly Lipschitz, so by Arzelà--Ascoli a subsequence converges to
  $\varphi:B(0,r)\to U\subset X$; by part (4) of
  \cref{prop:pet-ch11-norm-properties}, $\varphi$ is a homeomorphism onto its
  image, making $X$ a topological manifold. Writing the metrics
  $\varphi_i^*g_i$ in the fixed coordinates $B(0,r)$, another Arzelà--Ascoli
  argument (using that all $\varphi_i^*g_i$ satisfy (n1), (n2)) gives
  convergence in $C^m$ to a limit $g_*$ on $B(0,r)$ also satisfying (n1),
  (n2); these possibly only Hölder-continuous local metrics define local
  distances that agree with the restriction of $|\cdot|$ to $U$, since both
  the Riemannian and metric structures converge to the same limit on
  $B(0,r)$. Finally, for two overlapping charts $\varphi,\psi:B(0,r)\to X$,
  the transition $\varphi^{-1}\circ\psi$ is locally Lipschitz for the
  Euclidean metrics and distance preserving for the pulled-back metrics from
  $X$; a distance-preserving map between $C^\alpha$ Riemannian metrics is
  $C^{1,\alpha/2}$ (Calabi--Hartman), which is enough to give $X$ a
  $C^{1}$ (in fact $C^{m,\alpha}$, via the limiting metric components)
  differentiable structure compatible with these charts.
\end{proof}

\begin{theorem}[Thm.\ 11.3.6 — Fundamental Theorem of Convergence Theory]
  \label{thm:pet-ch11-fundamental-theorem-convergence}
  \lean{PetersenLib.fundamentalTheoremOfConvergenceTheory}
  \uses{lem:pet-ch11-fundamental-thm-precompact-gh,lem:pet-ch11-fundamental-thm-limit-is-manifold}
  For $Q>0$, $n\ge2$, $m\ge0$, $\alpha\in(0,1]$, $r>0$, the class
  $\mathcal{M}^{m,\alpha}(n,Q,r)$ of complete pointed Riemannian $n$-manifolds
  $(M,g,p)$ with $\|(M,g)\|_{C^{m,\alpha},r}\le Q$ is compact in the pointed
  $C^{m,\beta}$ topology for all $\beta<\alpha$.
  \source{petersen:page-0434,petersen:page-0435,petersen:page-0436,petersen:page-0437}
\end{theorem}
\begin{proof}
  \uses{thm:pet-ch11-fundamental-theorem-convergence}
  By \cref{lem:pet-ch11-fundamental-thm-precompact-gh}, a sequence
  $(M_i,g_i,p_i)\in\mathcal{M}$ Gromov--Hausdorff subconverges to some
  $(X,|\cdot|,p)$, which by \cref{lem:pet-ch11-fundamental-thm-limit-is-manifold}
  is a $C^{m,\alpha}$ Riemannian manifold in the same class. It remains to
  upgrade this to pointed $C^{m,\beta}$ convergence. Equip $X$ with a
  countable atlas $\varphi_s:B(0,r)\to U_s$ that are limits of charts
  $\varphi_{is}:B(0,r)\to U_{is}\subset M_i$ forming atlases of the $M_i$,
  with convergent transitions $\varphi_{is}^{-1}\circ\varphi_{jt}\to\varphi_s^{-1}\circ\varphi_t$
  and convergent metric components. Set
  $f_{is}=\varphi_{is}\circ\varphi_s^{-1}:U_s\to U_{is}$; then
  $(f_{is})^*g_i|_{U_s}\to g|_{U_s}$ in $C^{m,\beta}$, and the $f_{is}$
  ``converge to each other'' across overlapping charts in $C^{m+1,\beta}$.
  Build maps $F_{i\ell}:\Omega_\ell=\bigcup_{s=1}^\ell U_s\to\Omega_{i\ell}=\bigcup_{s=1}^\ell U_{is}$
  by induction on $\ell$: $F_{i1}=f_{i1}$; given $F_{i\ell}$, if
  $U_{\ell+1}\cap\Omega_\ell=\emptyset$ set $F_{i,\ell+1}=f_{i,\ell+1}$ on
  $U_{\ell+1}$ and $F_{i,\ell+1}=F_{i\ell}$ elsewhere; otherwise average, on
  $U_{\ell+1}$, via a partition of unity $\mu_1,\mu_2$ for $U_{\ell+1},\Omega_\ell$:
  \[
    F_{i,\ell+1}(x)=\varphi_{i,\ell+1}\circ\big(\mu_1(x)\varphi_{i,\ell+1}^{-1}\circ f_{i,\ell+1}(x)
    +\mu_2(x)\varphi_{i,\ell+1}^{-1}\circ F_{i\ell}(x)\big).
  \]
  In the coordinates $\varphi_{i,\ell+1}^{-1}\circ F_{i,\ell+1}\circ\varphi_{\ell+1}=\bar\mu_1F_1+\bar\mu_2F_2$,
  and since $\bar\mu_1F_1+\bar\mu_2F_2-F_1=\bar\mu_2(F_2-F_1)$ on all of
  $B(0,r)$ (the terms with $F_2$ undefined cancel against $\bar\mu_2=0$
  there),
  \[
    \|\bar\mu_1F_1+\bar\mu_2F_2-F_1\|_{C^{m+1,\beta}}\le C(n,m)\|\bar\mu_2\|_{C^{m+1,\beta}}\|F_2-F_1\|_{C^{m+1,\beta}}\to0
  \]
  as $i\to\infty$, since $F_2=\varphi_{i,\ell+1}^{-1}\circ F_{i\ell}\circ\varphi_{\ell+1}$
  and $F_1=\varphi_{i,\ell+1}^{-1}\circ f_{i,\ell+1}\circ\varphi_{\ell+1}$
  both converge to the identity-type limit transitions by the inductive
  hypothesis. Hence $F_{i,\ell+1}$ is $C^{m+1,\beta}$-close to $f_{is}$ for
  $s=1,\dots,\ell+1$ as $i\to\infty$, and closeness to the coordinate charts
  makes $F_{i\ell}$ an embedding on compact subsets of the domain. Letting
  $\ell\to\infty$ (with $i$ depending on $\ell$) exhibits embeddings
  $F_i:\Omega\to M_i$ realizing $(M_i,g_i,p_i)\to(X,g,p)$ in the pointed
  $C^{m,\beta}$ topology for every $\beta<\alpha$.
\end{proof}

\begin{corollary}[Cor.\ 11.3.7 — diameter/volume bounds give finiteness]
  \label{cor:pet-ch11-fundamental-thm-diam-vol-bounded}
  \lean{PetersenLib.fundamentalTheorem_diamOrVolBounded_finiteDiffeoTypes}
  \uses{thm:pet-ch11-fundamental-theorem-convergence}
  Subclasses of $\mathcal{M}^{m,\alpha}(n,Q,r)$ whose elements additionally
  satisfy $\diam\le D$, respectively $\vol\le V$, are compact in the
  $C^{m,\beta}$ topology (for $\beta<\alpha$); in particular each such
  subclass contains only finitely many diffeomorphism types.
  \source{petersen:page-0437}
\end{corollary}
\begin{proof}
  \uses{cor:pet-ch11-fundamental-thm-diam-vol-bounded}
  With notation as in the proof of
  \cref{thm:pet-ch11-fundamental-theorem-convergence}: if $\diam(M,g,p)\le D$
  then $M\subset B(p,k\delta/2)$ for $k=2D/\delta$, so each element of
  $\mathcal{M}^{m,\alpha}(n,Q,r)$ is covered by $\le N^k$ charts, and
  $C^{m,\beta}$ convergence is actually unpointed. If instead $\vol M\le V$,
  part (A) of the fundamental theorem gives that there are never more than
  $k=V\cdot e^{2nQ}\cdot(\vol B(0,\varepsilon))^{-1}$ disjoint
  $\varepsilon$-balls, so $\diam\le2\varepsilon k$ and the diameter case
  applies. Compactness in any $C^{m,\beta}$ topology implies only finitely
  many diffeomorphism types can occur.
\end{proof}

\begin{corollary}[Cor.\ 11.3.8 — compactness for the harmonic norm]
  \label{cor:pet-ch11-harmonic-norm-class-compact}
  \lean{PetersenLib.harmonicNormClass_closed_compact}
  \uses{prop:pet-ch11-harmonic-norm-properties,thm:pet-ch11-fundamental-theorem-convergence}
  For $Q>0$, $n\ge2$, $m\ge0$, $\alpha\in(0,1]$, $r>0$, the class of complete
  pointed Riemannian $n$-manifolds $(M,g,p)$ with
  $\|(M,g)\|^{\mathrm{har}}_{C^{m,\alpha},r}\le Q$ is closed in the pointed
  $C^{m,\alpha}$ topology and compact in the pointed $C^{m,\beta}$ topology
  for all $\beta<\alpha$.
  \source{petersen:page-0437}
\end{corollary}
\begin{proof}
  \uses{cor:pet-ch11-harmonic-norm-class-compact}
  By \cref{prop:pet-ch11-harmonic-norm-properties} the harmonic norm behaves
  like the ordinary norm under convergence, so
  \cref{thm:pet-ch11-fundamental-theorem-convergence} applies to give
  compactness; closedness follows because harmonic charts converge to
  harmonic charts, as established in the proof of
  \cref{prop:pet-ch11-harmonic-norm-properties}, so the limit space also has
  harmonic norm $\le Q$.
\end{proof}

\subsection{Alternative Norms}

\begin{remark}[alternative norm conditions]
  \label{rem:pet-ch11-alternative-norm-conditions}
  \lean{PetersenLib.alternativeNormConditions}
  \uses{def:pet-ch11-cmalpha-norm-pointed}
  The norm concept and its properties are unchanged if (n1), (n2) are
  replaced by
  \[
    \text{(n1$'$)}\quad|D\varphi|,|D\varphi^{-1}|\le f_1(n,Q),\qquad
    \text{(n2$'$)}\quad r^{|I|+\alpha}\|\partial^Ig\|_\alpha\le f_2(n,Q),\ 0\le|I|\le m,
  \]
  for continuous $f_1,f_2$ with $f_1(n,0)=1$, $f_2(n,0)=0$: this preserves
  continuity of the norm in $r$, the Fundamental Theorem, and the
  characterization of flat manifolds and Euclidean space.
  \source{petersen:page-0437}
\end{remark}

\begin{remark}[the $C^0$ norm and topological finiteness]
  \label{rem:pet-ch11-c0-norm-precompactness-only}
  \lean{PetersenLib.c0NormClass_precompact_gh_only}
  \uses{def:pet-ch11-cmalpha-norm-global,cor:pet-ch11-fundamental-thm-diam-vol-bounded}
  Taking $m=\alpha=0$, (n2) no longer makes sense but a $C^0$-norm concept
  survives. The class $\mathcal{M}^0(n,Q,r)$ is only precompact in the
  pointed Gromov--Hausdorff topology (the flat-manifold characterization
  still holds), and diameter- or volume-bounded subclasses are only
  Gromov--Hausdorff precompact, with diffeomorphism-type finiteness
  apparently failing: the averaging construction of
  \cref{thm:pet-ch11-fundamental-theorem-convergence} breaks down since the
  coordinates converge only in $C^{\alpha}$, $\alpha<1$. Using local
  contractibility of homeomorphism groups, two $C^0$-close topological
  embeddings can nonetheless be glued without much $C^0$ distortion, giving
  topological embeddings with converging pullback metrics; hence
  diameter- or volume-bounded classes contain only finitely many
  homeomorphism types, matching the original form of Cheeger's finiteness
  theorem (which considered diffeomorphism types under stronger coordinate
  bounds).
  \source{petersen:page-0438}
\end{remark}

\begin{definition}[Reifenberg norm]
  \label{def:pet-ch11-reifenberg-norm}
  \lean{PetersenLib.reifenbergNorm}
  \uses{def:pet-ch11-gromov-hausdorff-distance}
  For a metric space $(X,|\cdot|)$, the \emph{$n$-dimensional Reifenberg norm
  at scale $r$} is
  \[
    \|(X,|\cdot|)\|^n_r=\frac1r\sup_{p\in X}\dGH\big(B(p,r),B(0,r)\big),
  \]
  where $B(0,R)\subset\mathbb{R}^n$; the factor $r^{-1}$ prevents small distances
  merely because $r$ is small. If $(X_i,|\cdot|_i)\to(X,|\cdot|)$ in the
  Gromov--Hausdorff topology then $\|(X_i,|\cdot|_i)\|^n_r\to\|(X,|\cdot|)\|^n_r$
  for fixed $n,r$.
  \source{petersen:page-0438}
\end{definition}

\begin{remark}[Reifenberg norm: vanishing and the Cheeger--Colding converse]
  \label{rem:pet-ch11-reifenberg-norm-properties}
  \lean{PetersenLib.reifenbergNorm_limit_zero,PetersenLib.cheegerColdingReifenbergConverse}
  \uses{def:pet-ch11-reifenberg-norm}
  For an $n$-dimensional Riemannian manifold,
  $\lim_{r\to0}\|(M,g)\|^n_r=0$. Cheeger and Colding proved a converse: there
  is $\varepsilon(n)>0$ such that if $\|(X,|\cdot|)\|^n_r\le\varepsilon(n)$
  for all small $r$, then $X$ is, in a weak sense, an $n$-dimensional
  Riemannian manifold; in particular, for small $r$, the $\alpha$-Hölder
  distance between $B(p,r)$ and $B(0,r)$ is then small.
  \source{petersen:page-0438,petersen:page-0439}
\end{remark}

\begin{definition}[$\alpha$-Hölder distance]
  \label{def:pet-ch11-holder-distance}
  \lean{PetersenLib.holderDistance}
  \uses{def:pet-ch11-holder-pseudonorm}
  The \emph{$\alpha$-Hölder distance} $d_\alpha(X,Y)$ between metric spaces is
  the infimum, over homeomorphisms $F:X\to Y$, of
  \[
    \log\max\Big\{\sup_{x_1\ne x_2}\frac{|F(x_1)F(x_2)|}{|x_1x_2|^\alpha},\
    \sup_{y_1\ne y_2}\frac{|F^{-1}(y_1)F^{-1}(y_2)|}{|y_1y_2|^\alpha}\Big\}.
  \]
  \source{petersen:page-0438,petersen:page-0439}
\end{definition}

\begin{remark}[Cheeger--Colding Hölder convergence and Colding's volume criterion]
  \label{rem:pet-ch11-cheeger-colding-holder-convergence}
  \lean{PetersenLib.cheegerColdingHolderConvergence,PetersenLib.coldingVolumeCriterion}
  \uses{def:pet-ch11-holder-distance,rem:pet-ch11-reifenberg-norm-properties}
  Cheeger and Colding show that if $(M_i,g_i)\to(X,|\cdot|)$ in the
  Gromov--Hausdorff topology with $\|(M_i,g_i)\|^n_r\le\varepsilon(n)$ for all
  $i$ and small $r$, then $(M_i,g_i)\to(X,|\cdot|)$ in the Hölder distance,
  and in particular all $M_i$ are diffeomorphic to $X$ for large $i$.
  Colding earlier showed that for $(M,g)$ with $\Ric\ge(n-1)k$,
  $\|(M,g)\|^n_r$ is small iff $\vol B(p,r)\ge(1-\delta)\vol B(0,r)$ for some
  small $\delta$; relative volume comparison shows the volume condition for
  one $r$ persists for all smaller $r$, so smallness of the Reifenberg norm
  for all small $r$ follows from the volume condition at a single $r$.
  \source{petersen:page-0439}
\end{remark}

\section{Geometric Applications}

To obtain better control on the norms one uses harmonic coordinates on balls
of an a priori size, controlled by curvature and injectivity-radius bounds;
this section develops the resulting geometric applications, following
Anderson's approach.

\subsection{Ricci Curvature}

\begin{lemma}[Lemma 11.4.1 (Anderson, 1990) — Ricci controls the harmonic $C^{1,\alpha}$ norm]
  \label{lem:pet-ch11-anderson-ricci-harmonic-c1alpha-bound}
  \lean{PetersenLib.andersonRicciHarmonicC1AlphaBound}
  \uses{def:pet-ch11-harmonic-norm,lem:pet-ch11-laplacian-ricci-harmonic-coords,thm:pet-ch11-schauder-estimates}
  Suppose $(M,g)$ has $|\Ric|\le\Lambda$. For $r_1<r_2$,
  $K\ge\|(M,g,p)\|^{\mathrm{har}}_{C^{1,r_2}}$, $\alpha\in(0,1)$ there is
  $C(n,\alpha,K,r_1,r_2,\Lambda)$ with
  $\|(M,g,p)\|^{\mathrm{har}}_{C^{1,\alpha},r_1}\le C$. If moreover $g$ is
  Einstein, $\Ric=kg$, then for each $m$ there is
  $C(n,\alpha,K,r_1,r_2,k,m)$ with
  $\|(M,g,p)\|^{\mathrm{har}}_{C^{m+1,\alpha},r_1}\le C$.
  \source{petersen:page-0439,petersen:page-0440}
\end{lemma}
\begin{proof}
  \uses{lem:pet-ch11-anderson-ricci-harmonic-c1alpha-bound}
  Work in fixed harmonic coordinates on $B(0,r_2)$, where
  $\Delta=g^{ij}\partial_i\partial_j$. The bound
  $\|(M,g,p)\|^{\mathrm{har}}_{C^{1,r_2}}\le K$ gives the ellipticity data
  needed to apply \cref{thm:pet-ch11-schauder-estimates} to $\Delta g_{ij}$:
  \[
    \|g_{ij}\|_{C^{1,\alpha},B(0,r_1)}\le C(n,\alpha,K,r_1,r_2)\big(\|\Delta g_{ij}\|_{C^0,B(0,r_2)}+\|g_{ij}\|_{C^0,B(0,r_2)}\big).
  \]
  By \cref{lem:pet-ch11-laplacian-ricci-harmonic-coords},
  $\Delta g_{ij}=-2\Ric_{ij}-2Q(g,\partial g)$, so
  $\|\Delta g_{ij}\|_{C^0,B(0,r_2)}\le2\Lambda\|g_{ij}\|_{C^0,B(0,r_2)}+\hat C\|g_{ij}\|_{C^1,B(0,r_2)}$
  (bounding the quadratic-in-$\partial g$ term $Q$ using the $C^1$ bound from
  $K$), giving
  \[
    \|g_{ij}\|_{C^{1,\alpha},B(0,r_1)}\le C(n,\alpha,K,r_1,r_2)(2\Lambda+\hat C+1)\|g_{ij}\|_{C^1,B(0,r_2)}.
  \]
  In the Einstein case $\Ric=kg$ gives a $C^{1,\alpha}$ bound on $\Ric$
  itself (from the just-established $C^{1,\alpha}$ bound on $g$), so
  $\Delta g_{ij}$ is bounded in $C^\alpha$, not merely $C^0$; applying
  \cref{thm:pet-ch11-schauder-estimates} again gives a $C^{2,\alpha}$ bound
  on $g_{ij}$, whence $\Delta g_{ij}$ is bounded in $C^{1,\alpha}$, giving a
  $C^{3,\alpha}$ bound, and so on; iterating produces
  $C^{m+1,\alpha}$ bounds for every $m$.
\end{proof}

\begin{corollary}[Cor.\ 11.4.2 — precompactness under bounded Ricci and harmonic norm]
  \label{cor:pet-ch11-ricci-bounded-precompact-c1alpha}
  \lean{PetersenLib.ricciAndHarmonicNormBounded_precompact}
  \uses{lem:pet-ch11-anderson-ricci-harmonic-c1alpha-bound,thm:pet-ch11-fundamental-theorem-convergence}
  For $n\ge2$, $Q,r,\Lambda\in(0,\infty)$, the class of Riemannian
  $n$-manifolds with $\|(M,g)\|^{\mathrm{har}}_{C^{1,r}}\le Q$,
  $|\Ric|\le\Lambda$ is precompact in the pointed $C^{1,\alpha}$ topology for
  every $\alpha\in(0,1)$; the subclass of Einstein manifolds is compact in
  the $C^{m,\alpha}$ topology for every $m\ge1$, $\alpha\in(0,1)$.
  \source{petersen:page-0441}
\end{corollary}

\begin{theorem}[Thm.\ 11.4.3 (Anderson, 1990) — injectivity radius controls the harmonic norm]
  \label{thm:pet-ch11-anderson-injectivity-radius-controls-norm}
  \lean{PetersenLib.andersonInjectivityRadiusControlsNorm}
  \uses{lem:pet-ch11-anderson-ricci-harmonic-c1alpha-bound,thm:pet-ch11-fundamental-theorem-convergence,rem:pet-ch11-einstein-harmonic-regularity-bootstrap}
  Given $n\ge2$, $\alpha\in(0,1)$, $\Lambda,R>0$, for each $Q>0$ there is
  $r(n,\alpha,\Lambda,R)>0$ such that any compact Riemannian $n$-manifold
  $(M,g)$ with $|\Ric|\le\Lambda$, $\inj\ge R$ satisfies
  $\|(M,g)\|^{\mathrm{har}}_{C^{1,\alpha},r}\le Q$.
  \source{petersen:page-0441}
\end{theorem}
\begin{proof}
  \uses{thm:pet-ch11-anderson-injectivity-radius-controls-norm}
  Suppose not: there is $Q>0$ and, for each $i$, $(M_i,g_i)$ with
  $|\Ric|\le\Lambda$, $\inj\ge R$, but
  $\|(M_i,g_i)\|^{\mathrm{har}}_{C^{1,\alpha},r}>Q$ for all $r$. By continuity
  and vanishing at $r\to0$ of the norm, choose $r_i\in(0,i^{-1})$ with
  $\|(M_i,g_i)\|^{\mathrm{har}}_{C^{1,\alpha},r_i}=Q$. Rescale
  $\tilde g_i=r_i^{-2}g_i$; then $|\Ric(\tilde g_i)|\le r_i\Lambda$,
  $\inj(\tilde g_i)\ge r_i^{-1}R$, and
  $\|(M_i,\tilde g_i)\|^{\mathrm{har}}_{C^{1,\alpha},1}=Q$. Choose
  $p_i\in M_i$ with $\|(M_i,\tilde g_i,p_i)\|^{\mathrm{har}}_{C^{1,\alpha},1}\in[Q/2,Q]$.

  By \cref{lem:pet-ch11-anderson-ricci-harmonic-c1alpha-bound}, bounded Ricci
  curvature of $(M_i,\tilde g_i)$ bounds the $C^{1,\gamma}$ norm for
  $\gamma\in(\alpha,1)$, so by
  \cref{thm:pet-ch11-fundamental-theorem-convergence} a subsequence
  $(M_i,\tilde g_i,p_i)\to(M,g,p)$ in the pointed $C^{1,\gamma}$ topology,
  and continuity of the norm gives
  $\|(M,g,p)\|^{\mathrm{har}}_{C^{1,\alpha},1}\in[Q/2,Q]$.

  We claim $(M,g)=(\mathbb{R}^n,g_{\mathrm{eu}})$, contradicting the previous
  display. In converging harmonic coordinates, the metric coefficients
  satisfy $\tfrac12\Delta g_{kl}+Q(g,\partial g)=-\Ric_{kl}$
  (\cref{lem:pet-ch11-laplacian-ricci-harmonic-coords}); since
  $|\Ric(\tilde g_i)|\le r_i^{-2}\Lambda\cdot r_i^2=r_i\Lambda\to0$ (in the
  rescaled harmonic coordinates the Ricci bound tends to $0$) and
  $\tilde g_{i,kl}\to g_{kl}$ in $C^{1,\alpha}$, the limit satisfies
  $\tfrac12\Delta g_{kl}+Q(g,\partial g)=0$ weakly, so by
  \cref{rem:pet-ch11-einstein-harmonic-regularity-bootstrap} the limit is a
  smooth Ricci-flat manifold. Since $\inj(M_i,\tilde g_i)\to\infty$ and the
  metrics converge in $C^{1,\alpha}$, every finite-length geodesic in the
  limit is a limit of segments in $(M_i,\tilde g_i)$ and hence is itself a
  segment, giving $\inj(M,g)=\infty$; the splitting theorem then forces
  $(M,g)$ to be Euclidean space.
\end{proof}

\begin{corollary}[Cor.\ 11.4.4 (Anderson, 1990) — diffeomorphism finiteness under Ricci, diameter, injectivity bounds]
  \label{cor:pet-ch11-anderson-diffeo-finiteness-injectivity}
  \lean{PetersenLib.andersonDiffeoFinitenessInjectivity}
  \uses{thm:pet-ch11-anderson-injectivity-radius-controls-norm,cor:pet-ch11-fundamental-thm-diam-vol-bounded}
  For $n\ge2$, $\Lambda,D,R>0$, the class of closed Riemannian $n$-manifolds
  with $|\Ric|\le\Lambda$, $\diam\le D$, $\inj\ge R$ is precompact in the
  $C^{1,\alpha}$ topology for every $\alpha\in(0,1)$, and in particular
  contains only finitely many diffeomorphism types.
  \source{petersen:page-0442,petersen:page-0443}
\end{corollary}

\subsection{Volume Pinching}

\begin{lemma}[Lemma 11.4.5 (Anderson, 1990) — near-maximal volume growth forces Euclidean space]
  \label{lem:pet-ch11-anderson-volume-pinching-euclidean}
  \lean{PetersenLib.andersonVolumePinchingEuclidean}
  \uses{thm:pet-ch11-fundamental-theorem-convergence,cor:pet-ch11-norm-zero-implies-flat}
  For each $n\ge2$ there is $\varepsilon(n)>0$ such that any complete Ricci
  flat $(M,g)$ with $\vol B(p,r)\ge(1-\varepsilon)\omega_nr^n$ for some
  $p\in M$ and some $r$ is isometric to Euclidean space.
  \source{petersen:page-0443,petersen:page-0444}
\end{lemma}
\begin{proof}
  \uses{lem:pet-ch11-anderson-volume-pinching-euclidean}
  On any complete manifold with $\Ric\ge0$, relative volume comparison shows
  $\vol B(p,r)\ge(1-\varepsilon)\omega_nr^n$ for a single $p,r$ is equivalent
  (as $r$ varies, and over $p$) to
  $\lim_{r\to\infty}\vol B(p,r)/\omega_nr^n\ge1-\varepsilon$, so the
  hypothesis is independent of $p,r$ and is preserved under rescaling
  $g\mapsto\lambda^2g$ (since $\Ric\ge0$ is scale invariant).

  Suppose the lemma fails: for each $i$ there is Ricci-flat $(M_i,g_i)$ with
  $\lim_{r\to\infty}\vol B(p_i,r)/\omega_nr^n\ge1-i^{-1}$ but
  $\|(M_i,g_i)\|_{C^{1,\alpha},r}\ne0$ for all $r>0$. Rescaling, arrange
  $(M_i,\bar g_i,q_i)$ Ricci flat with
  $\lim_{r\to\infty}\vol B(q_i,r)/\omega_nr^n\ge1-i^{-1}$,
  $\|(M_i,\bar g_i)\|_{C^{1,\alpha},1}\le1$, and
  $\|(M_i,\bar g_i,q_i)\|_{C^{1,\alpha},1}\in[0.5,1]$. By
  \cref{thm:pet-ch11-fundamental-theorem-convergence} a subsequence converges
  in $C^{m,\alpha}$ to a Ricci-flat $(M,g,q)$; metric balls and volume forms
  converge, so $\lim_{r\to\infty}\vol B(q,r)/\omega_nr^n=1$, i.e.\ the limit
  has maximal volume growth for all balls, forcing $(M,g)$ to be Euclidean by
  the equality case of relative volume comparison
  (\cref{cor:pet-ch11-norm-zero-implies-flat} then applies once the norm is
  known to vanish). This contradicts continuity of the norm in $C^{1,\alpha}$
  (the norm of the limit would have to vanish, while
  $\|(M_i,\bar g_i,q_i)\|_{C^{1,\alpha},1}\ge0.5$ persists in the limit).
\end{proof}

\begin{corollary}[Cor.\ 11.4.6 — volume-pinching precompactness]
  \label{cor:pet-ch11-volume-pinching-precompactness}
  \lean{PetersenLib.volumePinchingPrecompactness}
  \uses{lem:pet-ch11-anderson-volume-pinching-euclidean,thm:pet-ch11-fundamental-theorem-convergence}
  For $n\ge2$, $-\infty<\lambda\le\Lambda<\infty$, $D,R\in(0,\infty)$, there
  is $\delta=\delta(n,\lambda R^2)$ such that the class of closed Riemannian
  $n$-manifolds with $(n-1)\Lambda\ge\Ric\ge(n-1)\lambda$, $\diam\le D$,
  $\vol B(p,R)\ge(1-\delta)v(n,\lambda,R)$ is precompact in the $C^{1,\alpha}$
  topology for every $\alpha\in(0,1)$, and in particular contains only
  finitely many diffeomorphism types.
  \source{petersen:page-0444,petersen:page-0445}
\end{corollary}
\begin{proof}
  \uses{cor:pet-ch11-volume-pinching-precompactness}
  As in the proof of \cref{thm:pet-ch11-anderson-injectivity-radius-controls-norm},
  suppose $(M_i,\bar g_i,p_i)$, $\bar g_i=k_i^2g_i$, $k_i\to\infty$, lie in
  the class. The volume hypothesis rescales to
  \[
    \vol_{\bar g_i}B(p_i,Rk_i)=k_i^n\vol_{g_i}B(p_i,R)\ge k_i^n(1-\delta)v(n,\lambda,R)
    =(1-\delta)v(n,\lambda k_i^{-2},Rk_i),
  \]
  and relative volume comparison then gives, for $r\le Rk_i$ and large $i$,
  $\vol_{\bar g_i}B(p_i,r)\ge(1-\delta)v(n,\lambda k_i^{-2},r)\sim(1-\delta)\omega_nr^n$.
  In the limit space (which exists by
  \cref{thm:pet-ch11-fundamental-theorem-convergence} applied to the
  precompact family obtained from bounded Ricci and diameter), this gives
  $\vol B(p,r)\ge(1-\delta)\omega_nr^n$ for all $r$; the limit is also Ricci
  flat (as $\lambda k_i^{-2}\to0$), so
  \cref{lem:pet-ch11-anderson-volume-pinching-euclidean} makes it Euclidean,
  contradicting continuity of the norm as before.
\end{proof}

\subsection{Sectional Curvature}

\begin{theorem}[Thm.\ 11.4.7 — the Convergence Theorem of Riemannian Geometry]
  \label{thm:pet-ch11-convergence-theorem-riemannian-geometry}
  \lean{PetersenLib.convergenceTheoremOfRiemannianGeometry}
  \uses{thm:pet-ch11-anderson-injectivity-radius-controls-norm}
  Given $R,K>0$ there exist $Q,r>0$ such that any $(M,g)$ with $\inj\ge R$,
  $|\sec|\le K$ has $\|(M,g)\|^{\mathrm{har}}_{C^{1,\alpha},r}\le Q$; in
  particular this class is compact in the pointed $C^{1,\alpha}$ topology for
  all $\alpha<1$.
  \source{petersen:page-0445}
\end{theorem}

\begin{corollary}[Cor.\ 11.4.8 (Cheeger, 1967) — even-dimensional pinched curvature finiteness]
  \label{cor:pet-ch11-cheeger-even-dim-finiteness}
  \lean{PetersenLib.cheegerEvenDimFiniteness}
  \uses{thm:pet-ch11-convergence-theorem-riemannian-geometry}
  For $n\ge1$, $k>0$, the class of Riemannian $2n$-manifolds with
  $k\le\sec\le1$ is compact in the $C^{1,\alpha}$ topology and hence contains
  only finitely many diffeomorphism types (using the diameter bound in
  positive curvature and Klingenberg's injectivity-radius estimate).
  \source{petersen:page-0445}
\end{corollary}

\begin{lemma}[Lemma 11.4.9 (Cheeger, 1967) — volume lower bound controls injectivity radius]
  \label{lem:pet-ch11-cheeger-injectivity-radius-lemma}
  \lean{PetersenLib.cheegerInjectivityRadiusLemma}
  \uses{thm:pet-ch11-fundamental-theorem-convergence,cor:pet-ch11-norm-zero-implies-flat}
  Given $n\ge2$, $v,K>0$, and a compact $n$-manifold $(M,g)$ with
  $|\sec|\le K$ and $\vol B(p,1)\ge v$ for all $p\in M$, then $\inj M\ge R$
  for $R=R(n,K,v)$.
  \source{petersen:page-0446}
\end{lemma}
\begin{proof}
  \uses{lem:pet-ch11-cheeger-injectivity-radius-lemma}
  Suppose not: $(M_i,g_i)$ satisfy the hypotheses with $\inj M_i\to0$. Choose
  $p_i$ with $\inj_{p_i}=\inj(M_i,g_i)$ and rescale
  $\bar g_i=(\inj M_i)^{-2}g_i$, so $\inj(M_i,\bar g_i)=1$ and
  $|\sec(M_i,\bar g_i)|\le(\inj M_i)^2K=:K_i\to0$. By
  \cref{thm:pet-ch11-fundamental-theorem-convergence}, a subsequence of
  $(M_i,\bar g_i,p_i)$ converges in the pointed $C^{1,\alpha}$ topology to
  $(M,g,p)$, flat since $\|(M,g)\|_{C^{1,\alpha},1}=0$.

  First, $\inj(M)\le1$: the conjugate radius for $(M_i,\bar g_i)$ is
  $\ge\pi/\sqrt{K_i}\to\infty$, so Klingenberg's estimate forces a geodesic
  loop of length $2$ at $p_i$; convergence carries this loop to a geodesic
  loop of length $2$ in $M$ based at $p$, so $\inj(M)\le1$.

  Second, $(M,g)=(\mathbb{R}^n,g_{\mathrm{std}})$: relative volume comparison and
  $\vol B(p_i,1)\ge v$ give $v'=v'(n,K,v)$ with
  $\vol B(p_i,r)\ge v'r^n$ for $r\le1$; the rescaled manifolds satisfy this
  for $r\le(\inj M_i)^{-1}$, and convergence gives $\vol B(p,r)\ge v'r^n$ for
  all $r$. Since $(M,g)$ is flat, this volume growth forces it to be
  Euclidean: if $M$ were a nontrivial quotient, the fundamental group would
  act by fixed-point-free isometries (infinite order, since finite-order
  isometries of Euclidean space are rotations about a fixed point), giving an
  intermediate cover $\mathbb{R}^n\to\tilde M\to M$ with $\pi_1(\tilde M)=\mathbb{Z}$,
  hence $\tilde M=\mathbb{R}^{n-1}\times S^1(R)$; then
  $\lim_{r\to\infty}\vol B(p,r)/r^{n-1}<\infty$ for $p\in\tilde M$ (and hence
  for $M$), contradicting $\vol B(p,r)\ge v'r^n$.

  The two conclusions ($\inj M\le1$ and $M$ Euclidean, with
  $\inj(\mathbb{R}^n,g_{\mathrm{std}})=\infty$) contradict each other.
\end{proof}

\begin{corollary}[Cor.\ 11.4.10 (Cheeger, 1967) — bounded curvature, diameter, volume finiteness]
  \label{cor:pet-ch11-cheeger-diam-vol-finiteness}
  \lean{PetersenLib.cheegerDiamVolFiniteness}
  \uses{lem:pet-ch11-cheeger-injectivity-radius-lemma,thm:pet-ch11-convergence-theorem-riemannian-geometry}
  For $n\ge2$, $K,D,v>0$, the class of closed Riemannian $n$-manifolds with
  $|\sec|\le K$, $\diam\le D$, $\vol\ge v$ is precompact in the $C^{1,\alpha}$
  topology for every $\alpha\in(0,1)$, and in particular contains only
  finitely many diffeomorphism types.
  \source{petersen:page-0447}
\end{corollary}

\subsection{Lower Curvature Bounds}

\begin{theorem}[Thm.\ 11.4.11 — lower sectional curvature and injectivity radius control the norm]
  \label{thm:pet-ch11-lower-sec-bound-injectivity-norm-control}
  \lean{PetersenLib.lowerSecBoundInjectivityNormControl}
  \uses{def:pet-ch11-cmalpha-norm-pointed}
  Given $R,k>0$ there exist $Q,r$ depending on $R,k$ such that any $(M,g)$
  with $\sec\ge-k^2$, $\inj\ge R$ satisfies $\|(M,g)\|_{C^{1,r}}\le Q$.
  \source{petersen:page-0447}
\end{theorem}
\begin{proof}
  \uses{thm:pet-ch11-lower-sec-bound-injectivity-norm-control}
  It suffices to bound the Hessian of distance functions $r(x)=|xp|$ on
  $B(p,R)-\{p\}$. The Hessian comparison theorem gives
  $\Hess r(x)\le k\coth(kr(x))g_r$. Conversely, for $r(x_0)<R$, $r(x)$ is
  supported from below near $x_0$ by $f(x)=R-|xx_0'|$, where $x_0'=c(R)$ and
  $c$ is the unit-speed minimizing geodesic from $p$ through $x_0$, giving
  $\Hess r\ge\Hess f\ge-k\coth(k(R-r(x_0)))g_r$ at $x_0$. Hence
  $|\Hess r|\le Q(k,R)$ on $B(x,r)$ whenever $|xp|\ge R/4$, $r\le R/4$.

  Fix $p\in M$, an orthonormal basis $e_1,\dots,e_n$ of $T_pM$, and geodesics
  $c_i$ with $c_i(0)=p$, $\dot c_i(0)=e_i$. Using
  $r^i(x)=|xc_i(-R/2)|$ on $B(p,R/4)$, define
  $\psi(x)=(r^1(x),\dots,r^n(x))-(r^1(p),\dots,r^n(p))$; then
  $D\psi|_p(e_i)$ is the standard basis of $T_0\mathbb{R}^n$, so $\psi$ is a
  chart near $p$ with $g_{ij}|_p=\delta_{ij}$. While $g_{ij}$ is not directly
  defined on all of $B(p,R/4)$, the (globally defined) potential inverse
  $g^{ij}=g(\nabla r^i,\nabla r^j)$ satisfies, by the Hessian estimate and
  $|\nabla r^k|=1$, $|dg^{ij}|\le Q(n,k,R)$ on $B(p,R/4)$, hence
  $\|g^{ij}-g^{ij}|_x\|<1/10$ for $x\in B(p,\delta(n,k,R))$. This shows
  $g^{ij}$ has a well-defined inverse $g_{ij}$ on $B(p,\delta)$ with
  $\|g_{ij}-g_{ij}|_x\|<1/9$ and $|dg_{ij}|\le C(n,k,R)$. By the inverse
  function theorem, $\psi$ is injective on $B(p,\delta)$ with
  $B(0,3/4)\subset\psi(B(p,\delta))$, and conditions (n1), (n2) hold with the
  resulting $Q,r$.
\end{proof}

\begin{example}[Example 11.4.12 — optimality: only $C^{1,1}$ limits under lower curvature]
  \label{ex:pet-ch11-optimality-lower-curvature-example}
  \lean{PetersenLib.optimalityLowerCurvatureExample}
  \uses{thm:pet-ch11-lower-sec-bound-injectivity-norm-control}
  For $\varepsilon\to0$, rotationally symmetric metrics
  $dr^2+\phi_\varepsilon^2(r)d\theta^2$ with $\phi_\varepsilon$ concave,
  $\phi_\varepsilon(r)=r$ for $0\le r\le1-\varepsilon$ and
  $\phi_\varepsilon(r)=\tfrac34r$ for $r\ge1+\varepsilon$, have $\sec\ge0$,
  $\inj\ge1$, and converge to a $C^{1,1}$ manifold with $C^{0,1}$ metric
  $(M,g)$, for which $\|(M,g)\|_{C^{1,r}}<\infty$ for all $r$: limit spaces
  of sequences with $\inj\ge R$, $\sec\ge-k^2$ cannot in general be smoother
  than this.
  \source{petersen:page-0448}
\end{example}

\begin{example}[Example 11.4.13 — a $C^{1,1}$ limit with two-sided curvature bound]
  \label{ex:pet-ch11-c11-limit-example}
  \lean{PetersenLib.c11LimitExample}
  \uses{ex:pet-ch11-optimality-lower-curvature-example}
  With $\psi_\varepsilon(r)=\sin r$ for $0\le r\le\pi/2-\varepsilon$ and
  $\psi_\varepsilon(r)=1$ for $r\ge\pi/2$, the metrics
  $dr^2+\psi_\varepsilon^2(r)d\theta^2$ satisfy $|\sec|\le4$, $\inj\ge1/4$; as
  $\varepsilon\to0$ they converge to a $C^{1,1}$ limit metric, though only
  $C^{1,\alpha}$ convergence for $\alpha<1$ has been established for such
  limits. Unlike the two-sided case, no injectivity-radius lower bound need
  hold under a one-sided curvature bound alone.
  \source{petersen:page-0448}
\end{example}

\begin{remark}[Exercise 11.4.14]
  \label{rem:pet-ch11-exercise-11-4-14}
  \lean{PetersenLib.exercise_11_4_14}
  \uses{ex:pet-ch11-optimality-lower-curvature-example}
  Given $a\in(0,1)$, $\varepsilon>0$, there is a smooth concave
  $\rho_\varepsilon(r)$ with $\rho_\varepsilon(r)=r$ for $0\le r\le\varepsilon$
  and $\rho_\varepsilon(r)=ar$ for $r\ge2\varepsilon$. The surfaces
  $dr^2+\rho_\varepsilon^2(r)d\theta^2$ have $\sec\ge0$, $\inj\le5\varepsilon$,
  and the volume of any $R$-ball is always $\ge a\pi R^2$.
  \source{petersen:page-0449}
\end{remark}

\begin{theorem}[Thm.\ 11.4.15 (Anderson--Cheeger, 1992) — Ricci lower bound with injectivity control]
  \label{thm:pet-ch11-anderson-cheeger-ricci-lower-bound-norm}
  \lean{PetersenLib.andersonCheegerRicciLowerBoundNorm}
  \uses{def:pet-ch11-harmonic-norm}
  Given $R,k>0$, $\alpha\in(0,1)$, there exist $Q,r$ depending on $n,R,k$
  such that any $(M^n,g)$ with $\Ric\ge-(n-1)k^2$, $\inj\ge R$ satisfies
  $\|(M,g)\|^{\mathrm{har}}_{C^{\alpha},r}\le Q$.
  \source{petersen:page-0449}
\end{theorem}

\subsection{Curvature Pinching}

\begin{theorem}[Thm.\ 11.4.16 — Ricci pinching implies $C^{1,\alpha}$-closeness to Einstein]
  \label{thm:pet-ch11-ricci-pinching-einstein-closeness}
  \lean{PetersenLib.ricciPinchingEinsteinCloseness}
  \uses{cor:pet-ch11-anderson-diffeo-finiteness-injectivity,rem:pet-ch11-einstein-harmonic-regularity-bootstrap}
  Given $n\ge2$, $R,D>0$, $\lambda\in\mathbb{R}$, there is
  $\varepsilon(n,\lambda,D,R)>0$ such that any closed Riemannian $n$-manifold
  $(M,g)$ with $\diam\le D$, $\inj\ge R$, $|\Ric-\lambda g|\le\varepsilon$ is
  $C^{1,\alpha}$ close to an Einstein metric with Einstein constant $\lambda$.
  \source{petersen:page-0450}
\end{theorem}
\begin{proof}
  \uses{thm:pet-ch11-ricci-pinching-einstein-closeness}
  By \cref{cor:pet-ch11-anderson-diffeo-finiteness-injectivity}, the class of
  such $(M,g)$ (for any fixed $\varepsilon$) is precompact in $C^{1,\alpha}$.
  If the theorem fails, there is a sequence $(M_i,g_i)\to(M,g)$ in
  $C^{1,\alpha}$ with $(M_i,g_i)$ never $C^{1,\alpha}$-close to an Einstein
  metric with constant $\lambda$, yet $|\Ric_{g_i}-\lambda g_i|\to0$. In
  harmonic coordinates, the limit metric is then a weak solution of
  $\tfrac12\Delta g+Q(g,\partial g)=-\lambda g$
  (\cref{lem:pet-ch11-laplacian-ricci-harmonic-coords}), so by
  \cref{rem:pet-ch11-einstein-harmonic-regularity-bootstrap} the limit is
  smooth and Einstein with constant $\lambda$ — contradicting that the
  $(M_i,g_i)$ were not close to such a metric.
\end{proof}

\begin{remark}[injectivity radius may be replaced by almost-maximal volume]
  \label{rem:pet-ch11-ricci-pinching-volume-alternative}
  \lean{PetersenLib.ricciPinchingVolumeAlternative}
  \uses{thm:pet-ch11-ricci-pinching-einstein-closeness,cor:pet-ch11-volume-pinching-precompactness}
  Using \cref{cor:pet-ch11-volume-pinching-precompactness} in place of the
  injectivity-radius precompactness in the proof of
  \cref{thm:pet-ch11-ricci-pinching-einstein-closeness}, the injectivity
  radius hypothesis there may be replaced by an almost-maximal-volume
  hypothesis.
  \source{petersen:page-0450}
\end{remark}

\begin{theorem}[Thm.\ 11.4.17 — sectional-curvature pinching implies closeness to constant curvature]
  \label{thm:pet-ch11-sec-pinching-constant-curvature-closeness}
  \lean{PetersenLib.secPinchingConstantCurvatureCloseness}
  \uses{thm:pet-ch11-ricci-pinching-einstein-closeness,lem:pet-ch11-cheeger-injectivity-radius-lemma}
  Given $n\ge2$, $v,D>0$, $\lambda\in\mathbb{R}$, there is
  $\varepsilon(n,\lambda,D,v)>0$ such that any closed Riemannian $n$-manifold
  $(M,g)$ with $\diam\le D$, $\vol\ge v$, $|\sec-\lambda|\le\varepsilon$ is
  $C^{1,\alpha}$ close to a metric of constant curvature $\lambda$.
  \source{petersen:page-0450}
\end{theorem}
\begin{proof}
  \uses{thm:pet-ch11-sec-pinching-constant-curvature-closeness}
  \cref{lem:pet-ch11-cheeger-injectivity-radius-lemma} gives a lower bound
  for $\inj$ from the curvature and volume bounds, so
  \cref{thm:pet-ch11-ricci-pinching-einstein-closeness} shows such metrics
  are $C^{1,\alpha}$ close to Einstein metrics with $\Ric=(n-1)\lambda g$. It
  remains to check that if $(M_i,g_i)\to(M,g)$ with
  $|\sec_{g_i}-\lambda|\to0$ and $\Ric_{g_i}=(n-1)\lambda g_i$, then $(M,g)$
  has constant curvature $\lambda$. Writing $g_i=dr^2+g_{r,i}$,
  $g=dr^2+g_r$ in polar coordinates about $p_i\to p$, convergence in
  $C^{1,\alpha}$ gives $g_{r,i}\to g_r$, and the Hessian comparison estimate
  for the curvature-pinched distance function gives
  \[
    \frac{\mathrm{sn}_{\lambda+\varepsilon_i}'(r_i)}{\mathrm{sn}_{\lambda+\varepsilon_i}(r_i)}g_{r,i}
    \le\Hess r_i\le\frac{\mathrm{sn}_{\lambda-\varepsilon_i}'(r_i)}{\mathrm{sn}_{\lambda-\varepsilon_i}(r_i)}g_{r,i}
  \]
  with $\varepsilon_i\to0$; passing to the limit,
  $\Hess r=\frac{\mathrm{sn}_\lambda'(r)}{\mathrm{sn}_\lambda(r)}g_r$, which
  forces the limit metric to have constant curvature $\lambda$ (this
  equality case is established elsewhere in the book).
\end{proof}

\begin{corollary}[Cor.\ 11.4.18 — even-dimensional positive pinching without hypotheses]
  \label{cor:pet-ch11-even-dim-positive-pinching}
  \lean{PetersenLib.evenDimPositivePinching}
  \uses{thm:pet-ch11-sec-pinching-constant-curvature-closeness}
  Given $2n\ge2$, $\lambda>0$, there is $\varepsilon=\varepsilon(n,\lambda)>0$
  such that any closed Riemannian $2n$-manifold $(M,g)$ with
  $|\sec-\lambda|\le\varepsilon$ is $C^{1,\alpha}$ close to a metric of
  constant curvature $\lambda$ (in even dimensions, positive curvature alone
  bounds the diameter and, together with an upper curvature bound, the
  injectivity radius from below).
  \source{petersen:page-0451}
\end{corollary}

\begin{corollary}[Cor.\ 11.4.19 — simply connected positive pinching]
  \label{cor:pet-ch11-odd-dim-simply-connected-pinching}
  \lean{PetersenLib.oddDimSimplyConnectedPinching}
  \uses{cor:pet-ch11-even-dim-positive-pinching}
  Given $n\ge2$, $\lambda>0$, there is $\varepsilon=\varepsilon(n,\lambda)>0$
  such that any closed simply connected Riemannian $n$-manifold $(M,g)$ with
  $|\sec-\lambda|\le\varepsilon$ is $C^{1,\alpha}$ close to a metric of
  constant curvature $\lambda$.
  \source{petersen:page-0451}
\end{corollary}

\section{Further Study}

Petersen closes the chapter with bibliographic remarks (no numbered
statements): Cheeger's original finiteness theorem and the $C^k$ convergence
theory for manifolds trace to his thesis and its journal version, with the
general pinching theorems also due to Cheeger; surveys by Anderson, Fukaya,
Petersen, and Yamaguchi, and by Petersen on norms and convergence, are
recommended for more on the theory and its applications, alongside the
French-language text that first popularized the Gromov--Hausdorff distance
and stronger convergence theorems (the present book states the Convergence
Theorem of Riemannian Geometry, \cref{thm:pet-ch11-convergence-theorem-riemannian-geometry},
in this form for the first time). S.\ Peters obtained an explicit estimate
for the number of diffeomorphism classes in Cheeger's finiteness theorem,
apparently the first place the modern statement of that theorem is proved.
Beyond quarter-pinching in positive curvature (proved earlier in the book
with purely topological conclusions), Ricci-flow techniques have since
sharpened the picture — first for positive curvature operator, then for
positive complex sectional curvature. In negative curvature, Heintze showed
that complete manifolds with $-1\le\sec<0$ have a lower volume bound once the
dimension is $\ge4$, making the volume hypothesis extraneous there, though
the diameter bound remains essential (with known counterexamples otherwise).
Near zero curvature, both the diameter and volume hypotheses are genuinely
needed — quotients of the Heisenberg group give bounded-diameter, arbitrarily
curvature-pinched examples admitting no flat metric — and Gromov classified
all $n$-manifolds with $|\sec|\le\varepsilon(n)$, $\diam\le1$ for small
$\varepsilon(n)$ as finitely covered by quotients of nilpotent Lie groups by
discrete subgroups, a result with an independent Ricci-flow proof as well.
\source{petersen:page-0451,petersen:page-0452}

\section{Exercises}

\begin{remark}[Exercise 11.6.1]
  \label{rem:pet-ch11-exercise-11-6-1}
  \lean{PetersenLib.exercise_11_6_1}
  \uses{def:pet-ch11-hausdorff-distance}
  Find a sequence of $1$-dimensional metric spaces that Hausdorff converge to
  the unit cube $[0,1]^3$ with the metric from the maximum norm on
  $\mathbb{R}^3$. Then find surfaces (jungle gyms) converging to the same
  space.
  \source{petersen:page-0453}
\end{remark}

\begin{remark}[Exercise 11.6.2]
  \label{rem:pet-ch11-exercise-11-6-2}
  \lean{PetersenLib.exercise_11_6_2}
  \uses{def:pet-ch11-gromov-hausdorff-distance}
  If $F:X\to Y$ (not necessarily continuous) satisfies
  $\big||x_1x_2|-|F(x_1)F(x_2)|\big|\le\varepsilon$ for all $x_1,x_2\in X$ and
  $F(X)\subset Y$ is $\varepsilon$-dense, show $\dGH(X,Y)<2\varepsilon$.
  \source{petersen:page-0453}
\end{remark}

\begin{remark}[Exercise 11.6.3]
  \label{rem:pet-ch11-exercise-11-6-3}
  \lean{PetersenLib.exercise_11_6_3}
  \uses{prop:pet-ch11-gromov-precompactness-criterion}
  (Croke) There is a universal constant $c(n)$ such that any $n$-manifold
  with $\inj\ge R$ satisfies $\vol B(p,r)\ge c(n)r^n$ for $r\le R/e$; use this
  to show the class of $n$-manifolds with $\inj\ge R$, $\vol\le V$ is
  precompact in the Gromov--Hausdorff topology.
  \source{petersen:page-0453}
\end{remark}

\begin{remark}[Exercise 11.6.4]
  \label{rem:pet-ch11-exercise-11-6-4}
  \lean{PetersenLib.exercise_11_6_4}
  \uses{cor:pet-ch11-ricci-diam-precompactness}
  Let $(M,g)$ be complete with $\Ric\ge(n-1)k$. Show there is
  $C(n,k)$ such that for each $\varepsilon\in(0,1)$ there is a cover by balls
  $B(x_i,\varepsilon)$ such that no more than $C(n,k)$ of the balls
  $B(x_i,5\varepsilon)$ have nonempty pairwise intersection.
  \source{petersen:page-0453}
\end{remark}

\begin{remark}[Exercise 11.6.5]
  \label{rem:pet-ch11-exercise-11-6-5}
  \lean{PetersenLib.exercise_11_6_5}
  \uses{lem:pet-ch11-laplacian-ricci-harmonic-coords}
  Show there are Bochner formulas for $\Hess(\tfrac12g(X,Y))$ and
  $\Delta g(X,Y)$, for vector fields $X,Y$ with symmetric $\nabla X,\nabla Y$,
  and use these to prove the formulas of
  \cref{lem:pet-ch11-laplacian-ricci-harmonic-coords} relating Ricci
  curvature to the metric in harmonic coordinates.
  \source{petersen:page-0453}
\end{remark}

\begin{remark}[Exercise 11.6.6]
  \label{rem:pet-ch11-exercise-11-6-6}
  \lean{PetersenLib.exercise_11_6_6}
  \uses{thm:pet-ch11-schauder-estimates}
  Show that, in contrast to the elliptic estimates of
  \cref{thm:pet-ch11-schauder-estimates}, it is not possible to find
  $C^{1,\gamma}$ bounds for a vector field $X$ in terms of $C^0$ bounds on
  $X$ and $\operatorname{div}X$.
  \source{petersen:page-0453}
\end{remark}

\begin{remark}[Exercise 11.6.7]
  \label{rem:pet-ch11-exercise-11-6-7}
  \lean{PetersenLib.exercise_11_6_7}
  \uses{def:pet-ch11-tensor-convergence-compact-subset,cor:pet-ch11-anderson-diffeo-finiteness-injectivity}
  Define $C^{m,\alpha}$ convergence for incomplete manifolds, with the
  boundary $\partial$ the points in the completion but not the manifold.
  Show that the class of incomplete spaces with $|\Ric|\le\Lambda$ and
  $\inj(p)\ge\min\{R,R\cdot d(p,\partial)\}$, $R<1$, is precompact in the
  $C^{1,\alpha}$ topology.
  \source{petersen:page-0453}
\end{remark}

\begin{remark}[Exercise 11.6.8]
  \label{rem:pet-ch11-exercise-11-6-8}
  \lean{PetersenLib.exercise_11_6_8}
  \uses{thm:pet-ch11-fundamental-theorem-convergence}
  Define a \emph{weighted norm}: fix $\rho(R)>0$ and require that, in a
  pointed manifold $(M,g,p)$, the points of $S(p,R)$ have norm $\le\rho(R)$.
  Prove the corresponding fundamental theorem.
  \source{petersen:page-0453}
\end{remark}

\begin{remark}[Exercise 11.6.9]
  \label{rem:pet-ch11-exercise-11-6-9}
  \lean{PetersenLib.exercise_11_6_9}
  \uses{def:pet-ch11-cmalpha-norm-global}
  If $\mathcal{M}$ is a class of compact Riemannian $n$-manifolds compact in
  the $C^{m,\alpha}$ topology, show there is $f(r)\to0$ as $r\to0$, depending
  on $\mathcal{M}$, with $\|(M,g)\|_{C^{m,\alpha},r}\le f(r)$ for all
  $M\in\mathcal{M}$.
  \source{petersen:page-0453}
\end{remark}

\begin{remark}[Exercise 11.6.10]
  \label{rem:pet-ch11-exercise-11-6-10}
  \lean{PetersenLib.exercise_11_6_10}
  \uses{def:pet-ch11-cmalpha-norm-global,thm:pet-ch11-fundamental-theorem-convergence}
  The \emph{local models} of a class of Riemannian manifolds are the limits
  obtained by rescaling elements of the class by constants $\to\infty$; e.g.\
  $|\sec|\le K$ rescales to $|\sec|\le\lambda^{-2}K\to0$, so the local models
  are the flat manifolds, and adding $\inj\ge R$ forces $\inj\to\infty$ under
  rescaling, so the only local model is Euclidean space. Show that a class of
  spaces is compact in the $C^{m,\alpha}$ topology if every sequence in the
  class, rescaled by constants $\to\infty$, subconverges in $C^{m,\alpha}$ to
  Euclidean space.
  \source{petersen:page-0454}
\end{remark}

\begin{remark}[Exercise 11.6.11]
  \label{rem:pet-ch11-exercise-11-6-11}
  \lean{PetersenLib.exercise_11_6_11}
  \uses{def:pet-ch11-gromov-hausdorff-distance}
  For the singular metric $dt^2+(at)^2d\theta^2$, $a>1$, on $\mathbb{R}^2$, show
  there is a sequence of rotationally symmetric metrics on $\mathbb{R}^2$ with
  $\sec\le0$, $\inj=\infty$ converging to it in the Gromov--Hausdorff
  topology.
  \source{petersen:page-0454}
\end{remark}

\begin{remark}[Exercise 11.6.12]
  \label{rem:pet-ch11-exercise-11-6-12}
  \lean{PetersenLib.exercise_11_6_12}
  \uses{thm:pet-ch11-fundamental-theorem-convergence}
  Show that the class of spaces with $\inj\ge R$ and
  $|\nabla^k\Ric|\le\Lambda$ for $k=0,\dots,m$ is compact in the
  $C^{m+1,\alpha}$ topology.
  \source{petersen:page-0454}
\end{remark}

\begin{remark}[Exercise 11.6.13]
  \label{rem:pet-ch11-exercise-11-6-13}
  \lean{PetersenLib.exercise_11_6_13}
  \uses{lem:pet-ch11-cheeger-injectivity-radius-lemma}
  (S.-h.\ Zhu) For the class of complete or compact $n$-manifolds with
  conjugate radius $\ge R$, $|\Ric|\le\Lambda$, $\vol B(p,1)\ge v$, use the
  techniques of \cref{lem:pet-ch11-cheeger-injectivity-radius-lemma} to show
  a lower injectivity-radius bound, hence compactness in $C^{1,\alpha}$.
  \source{petersen:page-0454}
\end{remark}

\begin{remark}[Exercise 11.6.14]
  \label{rem:pet-ch11-exercise-11-6-14}
  \lean{PetersenLib.exercise_11_6_14}
  \uses{rem:pet-ch11-exercise-11-6-13}
  Using the Eguchi--Hanson metrics, show that no compactness result holds in
  general for the class $|\Ric|\le\Lambda$, $\vol B(p,1)\ge v$ alone; either a
  large lower volume bound (as before) or a lower conjugate-radius bound must
  be assumed.
  \source{petersen:page-0454}
\end{remark}

\begin{remark}[Exercise 11.6.15]
  \label{rem:pet-ch11-exercise-11-6-15}
  \lean{PetersenLib.exercise_11_6_15}
  \uses{def:pet-ch11-harmonic-norm}
  The \emph{weak (harmonic) norm} $\|(M,g)\|^{\mathrm{weak}}_{C^{m,\alpha},r}$
  is defined as in \cref{def:pet-ch11-harmonic-norm}, except the charts
  $\varphi_s:B(0,r)\to U_s$ are only required to be immersions, so the
  inverse is only locally defined but still makes sense to be harmonic.
  \begin{enumerate}
  \item Show that if $(M,g)$ has bounded sectional curvature, then for every
    $Q>0$ there is $r>0$ with $\|(M,g)\|^{\mathrm{weak}}_{C^{m,\alpha},r}\le Q$;
    thus the weak norm is a generalized curvature quantity.
  \item Show the class of manifolds with bounded weak norm is precompact in
    the Gromov--Hausdorff topology.
  \item Show $(M,g)$ is flat iff the weak norm vanishes on all scales.
  \end{enumerate}
  \source{petersen:page-0454,petersen:page-0455}
\end{remark}
